\documentclass[11pt]{article}

\usepackage[T1]{fontenc}
\usepackage{lmodern}
\usepackage{amsmath,amssymb,amsthm}
\usepackage{mathtools}
\usepackage{microtype}
\usepackage{graphicx}
\usepackage{booktabs}
\usepackage{xcolor}
\usepackage{fancyhdr}
\usepackage[margin=1in]{geometry}
\usepackage[section]{placeins}   

\renewcommand{\topfraction}{0.85}
\renewcommand{\bottomfraction}{0.85}
\renewcommand{\textfraction}{0.10}
\renewcommand{\floatpagefraction}{0.70}

\makeatletter
\setlength{\@fptop}{0pt}
\setlength{\@fpbot}{0pt plus 1fil}
\makeatother
\usepackage{hyperref}
\hypersetup{
  colorlinks=true,
  linkcolor=blue!70!black,
  citecolor=green!50!black,
  urlcolor=blue!80!black,
  bookmarksnumbered=true,
  pdfauthor={Ioannis Tsiokos},
  pdftitle={Persistence of SPT-Legality from Channelized Feasibility in a Multiscale Hilbert/Cell Model}
}
\usepackage[numbers]{natbib}
\usepackage{orcidlink}

\let\H\relax
\newcommand{\H}{\mathcal{H}}
\newcommand{\C}{\mathcal{C}}
\newcommand{\F}{\mathcal{F}}
\newcommand{\conc}{\mathrm{conc}}
\newcommand{\kappaFun}{\kappa}
\newcommand{\proj}{P}
\newcommand{\Ran}{\mathrm{Ran}}
\newcommand{\Id}{\mathrm{Id}}

\newcommand{\R}{\mathbb{R}}
\newcommand{\norm}[1]{\left\lVert #1 \right\rVert}
\newcommand{\abs}[1]{\left\lvert #1 \right\rvert}
 
\newtheorem{definition}{Definition}[section]
\newtheorem{theorem}[definition]{Theorem}
\newtheorem{proposition}[definition]{Proposition}
\newtheorem{remark}[definition]{Remark}
\newtheorem{example}[definition]{Example}

\widowpenalty=10000
\clubpenalty=10000
\raggedbottom

\title{Persistence of SPT-Legality from Channelized Feasibility\\
in a Multiscale Hilbert/Cell Model}
\author{Ioannis Tsiokos\,\orcidlink{0009-0009-7659-5964}}
\date{15 February 2026}

\newcommand{\printaffiliation}{\begin{center}\small \end{center}}

\fancypagestyle{firstpage}{\fancyhf{}
  \fancyfoot[C]{\scriptsize Automorph Inc., Wilmington, DE, USA\;\texttt{ioannis@automorph.io}\;DOI: TBD\;\textcopyright\ 2026\;CC-BY 4.0\;Preprint --- not peer reviewed.}
  \renewcommand{\headrulewidth}{0pt}
  \renewcommand{\footrulewidth}{0pt}
}

\begin{document}
\maketitle
\thispagestyle{firstpage}
\printaffiliation

\begin{abstract}
Certifying multiscale anti-localization only at $t=0$ is not useful unless legality is forward-closed; we lack general, mechanizable conditions under which SPT-legality propagates under time evolution.
We propose a slot-audited channelized-feasibility contract in a Hilbert/band/cell model, where legality is defined by the normalized max-cell concentration audit $\kappa_j$ and enforced structurally by feasible representations $\mathcal{F}_j=\mathrm{Ran}(\Gamma_j)$ with packaging projectors $E_j$.
We introduce these objects and theorem-grade constants (delocalization, frame bounds, mode growth, and cellwise defect), operationalize them with standardized experiment outputs and an automated \texttt{audit.json}, and evaluate the resulting bounds in toy dynamics together with matched failure-mode controls.
In an invariant-ladder regime (exact feasibility invariance), $\kappa_j(t)$ stays uniformly below the predicted envelope $K^{\mathrm{pred}}_{\mathrm{glob}}$; injecting a globally tiny but cell-localized off-feasible component causes immediate deep-band violations despite leakage below $10^{-2}$.
With per-channel delocalization held fixed, worsening frame conditioning produces extreme $\kappa$ tails, whereas well-conditioned artifacts keep $\kappa$ far below the envelope; under cellwise defect control $\mathrm{def}^{\infty}\le \varepsilon_\infty f_j$, maxima grow monotonically with $\varepsilon_\infty$ and remain within the robust bound $K^{\mathrm{rob}}(\varepsilon_\infty)$.
This yields a closed, regression-tested regime for ``SPT-legal evolution'' in the model and pinpoints which structural slots are necessary (and which global diagnostics are misleading) when building auditable multiscale representations.
We are explicit that stronger claims---PDE smoothing, physical entropy production, or emergent learning of channel maps---are not established here; our conclusions hold in a finite-dimensional Hilbert/cell setting and rely on theorem-aligned bounds together with controlled counterexamples.
\end{abstract}

\noindent\textbf{keywords:} SPT-legality; channelized feasibility; anti-localization; multiscale; concentration audit; Hilbert/cell model; Six Birds framework

\section{Introduction}\label{sec:intro}

\subsection{Legality as a closed regime}

A recurring pattern in multiscale modeling is that one certifies a state as ``well-behaved'' at time $t=0$---for example, spread-out rather than wavepacket-like at fine scales---and then proceeds as if this property holds for the subsequent evolution.
From a mathematical perspective this is unsatisfactory unless the property defines a \emph{closed regime}: once the evolution enters the regime, it cannot leave without violating explicit structural hypotheses.
The goal of this paper is to formulate and prove such a closure statement for an anti-localization notion we call \emph{SPT-legality}.

We deliberately avoid any PDE-specific smoothing theory.
Instead we work in a minimal finite-dimensional setting and express hypotheses in the ``structural'' language promoted by the Six-Birds foundations paper \cite{Tsiokos2026SixBirds}:
packaging (idempotent representations), gating/feasibility constraints, staged interface growth, and explicit audits.
In this way, ``legality persistence'' becomes a theorem about compatibility between dynamics and representation, rather than a theorem about analytic regularization.

\subsection{A minimal multiscale Hilbert/cell model}

Let $H$ be a real Hilbert space with an orthogonal multiscale decomposition
\[
H \;=\; \bigoplus_{j=0}^{J_{\max}} H_j,
\qquad u_j := P_j u \in H_j,
\]
and, for each depth $j$, a fixed orthogonal ``cell'' partition
\[
H_j \;\cong\; \bigoplus_{c\in \mathcal C_j} H_{j,c},
\qquad P_{j,c}:H_j\to H_{j,c}.
\]
Write the baseline cell fraction $f_j := 1/|\mathcal C_j|$.
For a nonzero $v\in H_j$, define the maximum cell energy fraction
\[
\mathrm{conc}_j(v) \;:=\; \max_{c\in\mathcal C_j}\frac{\|P_{j,c}v\|^2}{\|v\|^2},
\]
and the normalized concentration
\[
\kappa_j(v) \;:=\; \frac{\mathrm{conc}_j(v)}{f_j}.
\]
Thus $\kappa_j(v)\approx 1$ indicates that $v$ is spread roughly uniformly across cells at depth $j$, while $\kappa_j(v)\gg 1$ indicates localization.

\paragraph{SPT-legality.}
Fix a threshold $K_\kappa\ge 1$.
A state $u\in H$ is \emph{SPT-legal} if for every depth $j$ with $u_j\neq 0$ one has $\kappa_j(u_j)\le K_\kappa$.
The persistence problem asks: under what conditions on the dynamics does SPT-legality at $t=0$ imply SPT-legality for all later times?

\subsection{Six-Birds framing: packaging, gating, staging, audits}

The Six-Birds foundations paper \cite{Tsiokos2026SixBirds} emphasizes that closed regimes arise from a small set of primitives:
\emph{packaging} (P5) provides an idempotent representation map whose fixed points define ``objects'';
\emph{gating/constraints} (P2) restrict which patterns are reachable;
\emph{sectors and staged interface growth} (P4) control complexity across refinement;
and \emph{audits/accounting} (P6) provide explicit certificates.
Route-mismatch/holonomy (P3) appears as a defect signal when packaging and evolution do not commute, and operator rewrite (P1) is required when one wants to restore closure by changing the representation.

Our contribution is to instantiate this program concretely.
At each depth $j$ we introduce a \emph{channel map} $\Gamma_j(a)=\sum_{r=1}^{m_j}a_r\psi_{j,r}$ with a feasible set $\mathcal F_j=\mathrm{Ran}(\Gamma_j)\subseteq H_j$.
Packaging is modeled by the orthogonal projector $E_j$ onto $\mathcal F_j$.
Delocalization and frame bounds are treated as \emph{artifact} properties of $\Gamma_j$ (auditable numerically), while invariance or controlled defect of feasibility is treated as the \emph{system} compatibility property that enables persistence.

\subsection{Main results (informal)}

Our results separate into three layers.

\paragraph{(i) Static feasibility implies legality bounds.}
If $v\in\mathcal F_j$ is representable through a channel map whose columns are cell-delocalized and whose Gram matrix has controlled condition number, then $\kappa_j(v)$ admits an explicit bound of the form
\[
\kappa_j(v)\;\lesssim\; \kappa_0(j)\,\mathrm{cond}(\Gamma_j)\,m_j,
\]
where $\kappa_0(j)$ is a delocalization constant and $m_j$ is the channel count at depth $j$.

\paragraph{(ii) Exact invariance implies persistence.}
If a dynamics preserves feasibility bandwise (i.e.\ $u_j(t)\in\mathcal F_j$ for all $t$), then the static bound propagates forward in time: SPT-legality is forward-closed with threshold $K=\max_j K^{\mathrm{feas}}_j$.

\paragraph{(iii) Robust persistence from cellwise defect control.}
Exact invariance is often too strong.
We therefore prove a robust variant: if $(I-E_j)u_j(t)$ is small in a \emph{cellwise} sense, then SPT-legality persists with an explicit degraded bound $K^{\mathrm{rob}}(\varepsilon)$.
A key point is that global distance-to-feasibility is not the right robustness notion; the relevant defect must be local in the cell partition.

\subsection{Sharpness: what must fail if a slot is dropped}

We also provide no-go results and counterexamples clarifying why persistence cannot be proved from accounting alone.

\begin{itemize}
\item \textbf{Ledger is not legality:} even with exactly conserved energy, $\kappa_j$ can approach its maximum $1/f_j$ along feasible directions if localized feasible patterns exist.
\item \textbf{Delocalized columns are not enough:} without a frame lower bound, near-dependence produces coefficient directions with extreme concentration tails.
\item \textbf{Tiny global leakage can hide deep illegality:} global norm leakage can be arbitrarily small while deep-band $\kappa_j$ becomes near-maximal.
\end{itemize}

These boundaries identify the true ``hard slots'' needed for a closed legal regime: structural anti-localization in feasibility, controlled channel count growth, and nondegenerate frame bounds, together with (exact or robust) invariance of feasibility.

\subsection{Numerical hardening and assumption audits}

To harden the theorem overlays and to support the sharpness claims, we include numerical experiments built directly in this Hilbert/cell model.
Each run emits machine-readable outputs and an automated assumption audit report (\texttt{audit.json}) that checks delocalization, frame conditioning, mode growth, invariance/defect evidence, and intended failure modes.
These experiments serve as regression tests against silent drift while the paper is written.

\subsection{Contributions}

\begin{itemize}
\item We formalize SPT-legality as a quantitative anti-localization condition in a finite-dimensional Hilbert/band/cell model, with a simple concentration audit functional $\kappa_j$.
\item We introduce channelized feasibility sets and packaging projectors as the structural mechanism enforcing anti-localization, in the spirit of the Six-Birds primitives \cite{Tsiokos2026SixBirds}.
\item We state and prove a persistence theorem for SPT-legality under exact feasibility invariance, and a robust variant under cellwise defect control, with explicit constants.
\item We prove sharp no-go boundaries showing that ledger invariants, delocalized columns without frame bounds, or global leakage metrics cannot guarantee legality persistence.
\item We provide toy models and numerical experiments (with automated assumption audits) that validate the happy-path regime and isolate each failure mode.
\end{itemize}

\subsection{Organization}

Section~2 restates the Six-Birds primitives (P1--P6) and maps them into the present Hilbert/cell setting.
Section~3 gives formal definitions of concentration and legality, channel maps and feasible sets, packaging projectors, and artifact diagnostics.
Sections~4--5 state the hypothesis slots and the main persistence theorems.
Sections~6--8 contain proofs.
Sections~9--10 present non-implications and no-go results.
Section~11 gives toy models, and Section~12 reports numerical experiments and audits.
We close with an assumption audit (Section~13), a conclusion (Section~14), and a reproducibility appendix (Appendix~A).
 \section{Relation to the Six-Birds primitives (P1--P6)}\label{sec:six-birds}

This paper is intentionally written in the structural style advocated by the Six-Birds foundations paper \cite{Tsiokos2026SixBirds}.
Rather than relying on PDE-specific regularization, we identify a multiscale ``legal regime'' (anti-localization across cells) and ask for conditions under which this regime is \emph{closed} under time evolution.
In the Six-Birds vocabulary, closure is not something one gets ``for free'' from accounting alone; it must be enforced by compatibility between dynamics and representation (packaging and gating), together with staged control of interface complexity and explicit audits.

\subsection{The primitives (informal restatement)}

The foundations paper \cite{Tsiokos2026SixBirds} proposes six primitives (``birds'') for emergence calculus.
We restate them informally here, in the minimal form needed for the present work.

\begin{description}
\item[\textbf{P1: Operator rewrite.}]
When closure fails, one often must change the effective operator/\hspace{0pt}representation used by the theory---not merely iterate the same dynamics.
This is the ``rewrite'' move: replace the map one is studying by another that is compatible with the intended packaged objects.

\item[\textbf{P2: Gating / constraints.}]
Not all patterns are reachable or admissible.
Gating restricts feasible macrostates to those representable through a finite interface, thereby eliminating unstable or non-representable degrees of freedom.

\item[\textbf{P3: Route mismatch / holonomy (diagnostic).}]
If two routes through a reduction pipeline do not commute (e.g.\ evolve then package vs.\ package then evolve), the mismatch provides a diagnostic signal.
Route mismatch is not itself a certificate of control; it quantifies incompatibility.

\item[\textbf{P4: Sectors and staged interface growth.}]
A system may decompose into sectors (invariant labels) and a staged notion of depth/refinement.
Crucially, the effective interface size (degrees of freedom needed to communicate with deeper scales) must remain controlled with depth for uniform statements to be meaningful.

\item[\textbf{P5: Packaging.}]
Packaging is an idempotent map whose fixed points define the ``objects'' of the theory.
A packaged description discards irrelevant detail while retaining the invariant representation needed for subsequent evolution.

\item[\textbf{P6: Accounting / audit.}]
Audits are explicit, checkable certificates (often monotone quantities) that verify regime membership or rule out certain failures.
Accounting alone does not force closure, but it can be used to \emph{detect} when closure hypotheses are violated.
\end{description}

\subsection{Instantiation in the Hilbert/band/cell model}

Our multiscale setting (Section~\ref{sec:intro} and formalized in Section~3) provides concrete realizations of each primitive:

\paragraph{P5 (packaging) as bandwise projectors onto feasible sets.}
At depth $j$ we introduce a synthesis/channel map $\Gamma_j:\mathbb{R}^{m_j}\to H_j$ and define the feasible set
\[
\mathcal{F}_j := \mathrm{Ran}(\Gamma_j)\subseteq H_j.
\]
Packaging is modeled by the orthogonal projector
\[
E_j: H_j \to \mathcal{F}_j,
\qquad E_j^2 = E_j,
\qquad \mathcal{F}_j = \{v\in H_j:\;E_j v = v\}.
\]
In this paper, ``feasible at depth $j$'' is synonymous with ``packaged at depth $j$.''

\paragraph{P2 (gating) as feasibility invariance or controlled feasibility defect.}
Gating enters as a compatibility condition between the dynamics and the packaged feasible sets:
either \emph{exact invariance} (for all times, $u_j(t)\in\mathcal{F}_j$) or \emph{robust invariance} (the deviation $(I-E_j)u_j(t)$ is controlled).
This is the core structural assumption required for legality persistence.

\paragraph{P4 (staging and interface growth) as controlled channel count $m_j$ across depth.}
The depth index $j$ is the staging parameter, and the interface size at depth $j$ is captured by the channel count $m_j$.
Uniform legality across depth is only plausible when $m_j$ grows slowly enough relative to the cell refinement $|\mathcal{C}_j|$.
The ``mode explosion'' failure mode in our experiments illustrates that if $m_j$ grows too fast (e.g.\ to full dimension), a fixed legality threshold cannot hold across depth.

\paragraph{P6 (audit) as concentration legality checks and artifact diagnostics.}
Our legal regime is defined by an explicit audit functional: the normalized concentration $\kappa_j$ based on the cell partition.
In addition, feasibility artifacts $\Gamma_j$ admit their own auditable diagnostics: empirical delocalization constants, frame bounds (via Gram spectra), and predicted feasibility legality envelopes.
In the codebase, every run emits an \texttt{audit.json} that checks these slots numerically (Section~12).

\paragraph{P3 (route mismatch) as a defect functional measured at cell resolution.}
When packaging and evolution do not commute, the relevant diagnostic is not merely a global distance to feasibility but a \emph{cellwise} mismatch.
Our robust persistence theorem is formulated in terms of cellwise defect measures of $(I-E_j)u_j(t)$, reflecting the fact that legality is itself a cell-level anti-localization notion.

\paragraph{P1 (operator rewrite) as the remedy when invariance fails.}
If the dynamics is not compatible with the packaged feasible sets (or if route mismatch is uncontrolled), legality cannot be forward-closed.
In the Six-Birds language, the correct response is not to hope that an accounting invariant will prevent failure, but to rewrite: change the representation/packaging (the maps $\Gamma_j$ or $E_j$) or modify the effective operator so that feasibility becomes (approximately) invariant and auditable.

\begin{table}[t]
\centering
\caption{Six-Birds primitives \cite{Tsiokos2026SixBirds} and their instantiation in the present Hilbert/band/cell legality framework.}
\label{tab:six-birds-mapping}
\begin{tabular}{@{}p{1.8cm}p{5.4cm}p{7.8cm}@{}}
\toprule
\textbf{Primitive} & \textbf{Role in Six-Birds} & \textbf{Instantiation here} \\
\midrule
\textbf{P1} & Operator rewrite when closure fails. & Modify packaging/feasibility maps $(\Gamma_j,E_j)$ or the effective evolution so feasibility becomes invariant/robustly invariant. \\
\textbf{P2} & Gating/constraints restrict reachable patterns. & Bandwise feasibility: $u_j(t)\in\mathcal{F}_j=\mathrm{Ran}(\Gamma_j)$ (exact) or controlled defect $(I-E_j)u_j(t)$ (robust). \\
\textbf{P3} & Route mismatch/holonomy is a diagnostic of incompatibility. & Cellwise defect metrics for $(I-E_j)u_j(t)$; ``evolve then package'' vs.\ ``package then evolve'' mismatch. \\
\textbf{P4} & Sectors and staged interface growth. & Depth $j$ as refinement stage; channel count $m_j$ as interface size; mode-growth control prevents uniform legality collapse. \\
\textbf{P5} & Packaging as an idempotent map; fixed points are objects. & $E_j$ orthogonal projector onto $\mathcal{F}_j$; packaged (feasible) band states satisfy $E_j v=v$. \\
\textbf{P6} & Accounting/audit certificates. & Legality audit $\kappa_j$ and artifact diagnostics (empirical delocalization, frame bounds); emitted per run as \texttt{audit.json}. \\
\bottomrule
\end{tabular}
\end{table}

\subsection{How this framing structures the paper}

With this mapping in place, our main theorems can be read as ``closure of a P6-defined legal regime'' under P5/P2/P4 compatibility hypotheses, with P3-style defect control providing robustness.
The no-go results then identify which primitive assumptions are genuinely necessary: accounting invariants do not prevent concentration blow-up if the feasible sets admit localized directions, and delocalized columns without frame control do not prevent concentrated combinations.
 \section{Model and Definitions}\label{sec:definitions}

This section formalizes the minimal multiscale setting used throughout the paper.
All statements are finite-dimensional (for mechanizability), but the definitions are written in Hilbert language.
The intended interpretation is: depth $j$ indexes refinement, and the cell partition at depth $j$ represents the resolution at which localization is measured.

\subsection{Multiscale bands and cell partitions}

Fix a real Hilbert space $H$ with an orthogonal decomposition into depth/band subspaces
\begin{equation}\label{eq:hilbert-decomp}
H \;=\; \bigoplus_{j=0}^{J_{\max}} H_j,
\qquad
P_j:H\to H_j
\text{ orthogonal projectors}.
\end{equation}
For a state $u\in H$ we write $u_j:=P_j u\in H_j$.

For each depth $j$ we also fix an orthogonal decomposition into ``cells''
\begin{equation}\label{eq:cell-decomp}
H_j \;\cong\; \bigoplus_{c\in \mathcal{C}_j} H_{j,c},
\qquad
P_{j,c}:H_j\to H_{j,c}
\text{ orthogonal projectors}.
\end{equation}
The cell count is $|\mathcal{C}_j|$ and we define the baseline cell fraction
\begin{equation}\label{eq:baseline-fj}
f_j \;:=\; \frac{1}{|\mathcal{C}_j|}.
\end{equation}

\subsection{Trajectories and lifespan}\label{subsec:traj}

A \emph{trajectory} is a map $u:I\to H$, where the time set $I$ is either an interval $I\subset\mathbb{R}$ (continuous time) or a discrete set $I\subset\mathbb{Z}$ (discrete time).
We refer to $I$ as the \emph{lifespan} of the trajectory.
All persistence statements in this paper are of the form: ``for all $t\in I$'' under explicit HL slots; we do not assume any PDE-specific regularity or smoothing beyond the existence of the trajectory in the stated time domain.

\subsection{Concentration and SPT-legality}

\begin{definition}[Cell concentration]\label{def:cell-conc}
For $j\le J_{\max}$ and nonzero $v\in H_j$, define the maximum cell energy fraction
\begin{equation}\label{eq:conc-def}
\conc_j(v) \;:=\; \max_{c\in\mathcal{C}_j}\frac{\|P_{j,c}v\|^2}{\|v\|^2}.
\end{equation}
We set $\conc_j(0):=0$ by convention.
\end{definition}

\begin{definition}[Normalized concentration]\label{def:kappa-def}
For $j\le J_{\max}$ and $v\in H_j$, define
\begin{equation}\label{eq:kappa-def}
\kappa_j(v) \;:=\; \frac{\conc_j(v)}{f_j}.
\end{equation}
Thus $\kappa_j(v)\in[0,1/f_j]$, with $\kappa_j(v)\approx 1$ corresponding to roughly uniform energy distribution across cells, while $\kappa_j(v)\approx 1/f_j$ corresponds to near-complete localization in a single cell.
\end{definition}

\begin{definition}[SPT-legal state]\label{def:spt-legal}
Fix a legality threshold $K_\kappa\ge 1$.
A state $u\in H$ is \emph{SPT-legal} (with threshold $K_\kappa$) if for every $j\le J_{\max}$ with $u_j\neq 0$ one has
\begin{equation}\label{eq:spt-legal}
\kappa_j(u_j)\le K_\kappa.
\end{equation}
\end{definition}

\subsection{Channel maps, feasible sets, and packaging}

We model ``structural anti-localization'' not by analytic smoothing but by representability constraints.

\begin{definition}[Channel map and feasible set]\label{def:channel-map}
Fix a channel count $m_j\in\mathbb{N}$ for each depth $j$.
A \emph{channel map} at depth $j$ is a linear synthesis operator
\begin{equation}\label{eq:gamma-def}
\Gamma_j:\mathbb{R}^{m_j}\to H_j,
\qquad
\Gamma_j a \;=\; \sum_{r=1}^{m_j} a_r\,\psi_{j,r},
\end{equation}
where $\psi_{j,r}\in H_j$ are the \emph{channels}.
The associated feasible set is the range
\begin{equation}\label{eq:feasible-set}
\mathcal{F}_j \;:=\; \Ran(\Gamma_j)\subseteq H_j.
\end{equation}
We call $\Gamma_j$ (equivalently, the channel family $\{\psi_{j,r}\}$) an \emph{artifact} at depth $j$.
\end{definition}

\begin{definition}[Packaging projector]\label{def:packaging}
The (bandwise) packaging map associated to $\mathcal{F}_j$ is the orthogonal projector
\begin{equation}\label{eq:packaging-proj}
E_j:H_j\to \mathcal{F}_j,
\qquad
E_j^2 = E_j,
\qquad
\mathcal{F}_j = \{v\in H_j:\;E_j v=v\}.
\end{equation}
We also write $I$ for the identity map on $H_j$ and refer to $(I-E_j)v$ as the \emph{off-feasible component} of $v$.
\end{definition}

\paragraph{Artifact vs.\ system.}
Artifacts $(\Gamma_j,E_j)$ are fixed data of a multiscale representation.
A dynamical system provides a trajectory $u(t)\in H$.
The persistence results later in the paper combine (i) \emph{artifact} properties (delocalization, frame bounds, mode growth) with (ii) \emph{system} compatibility assumptions (invariance or controlled defect relative to feasibility).

\subsection{Delocalization and frame bounds (artifact hypotheses)}

We will formulate theorems in terms of auditable constants associated to an artifact.

\begin{definition}[Channel delocalization]\label{def:deloc}
A channel family $\{\psi_{j,r}\}_{r=1}^{m_j}$ is \emph{cell-delocalized} with constant $\kappa_0(j)\ge 1$ if
\begin{equation}\label{eq:deloc-hyp}
\max_{c\in\mathcal{C}_j}\|P_{j,c}\psi_{j,r}\|^2
\;\le\;
\kappa_0(j)\, f_j\,\|\psi_{j,r}\|^2
\qquad
\text{for all } r.
\end{equation}
\end{definition}

\begin{definition}[Frame bounds and conditioning]\label{def:frame-bounds}
A channel map $\Gamma_j$ has (Euclidean) frame bounds $0<\alpha_j\le \beta_j<\infty$ if
\begin{equation}\label{eq:frame-bounds}
\alpha_j\,\|a\|^2 \;\le\; \|\Gamma_j a\|^2 \;\le\; \beta_j\,\|a\|^2
\qquad
\text{for all } a\in\mathbb{R}^{m_j}.
\end{equation}
We define the associated condition number proxy
\begin{equation}\label{eq:cond-def}
\mathrm{cond}(\Gamma_j) \;:=\; \beta_j/\alpha_j.
\end{equation}
In finite dimensions, one may take $\alpha_j,\beta_j$ as the smallest/largest eigenvalues of the Gram matrix $\Gamma_j^\top\Gamma_j$ when $\Gamma_j$ is full rank.
See, e.g., \cite{Christensen2016Frames,DuffinSchaeffer1952Nonharmonic} for background on frames/Riesz sequences and the role of frame bounds.
\end{definition}

\begin{remark}[Terminology and scope: stable synthesis vs.\ redundant frames]\label{rem:stable-synthesis}
The inequality \eqref{eq:frame-bounds} is a stability condition for the \emph{synthesis} map $a\mapsto \Gamma_j a$.
In classical frame theory, this corresponds to the channels forming a \emph{Riesz sequence} (equivalently, $\Gamma_j$ is injective and coefficients are stable), rather than an arbitrary redundant frame.
In particular, $\alpha_j>0$ forces $m_j\le \dim(H_j)$.
We adopt this strong form because our legality envelopes are stated in terms of the coefficient norm and are representation-independent only when coefficients are unique (or otherwise canonically selected).
Extending the theory to redundant dictionaries would require fixing a coefficient-selection rule (e.g.\ minimal $\ell_2$-norm coefficients) and replacing $\alpha_j$ by a pseudoinverse stability constant; we do not pursue this here.
\end{remark}

\subsection{Artifact diagnostics used in theorem overlays and audits}

The experiments and automated audits compute empirical versions of the above constants.

\begin{definition}[Empirical delocalization constant]\label{def:kappa0-emp}
Given channels $\{\psi_{j,r}\}$ at depth $j$, define
\begin{equation}\label{eq:kappa0-emp}
\kappa_0^{\mathrm{emp}}(j)
\;:=\;
\max_{r}\max_{c\in\mathcal{C}_j}
\frac{\|P_{j,c}\psi_{j,r}\|^2}{f_j\,\|\psi_{j,r}\|^2},
\end{equation}
In our finite-dimensional setting this is an exact worst-case artifact constant computed from the channels; the superscript ``emp'' emphasizes that it is audited from artifact data rather than estimated from trajectory samples.
with the convention that zero-norm channels are ignored.
\end{definition}

\begin{definition}[Predicted feasibility legality envelope]\label{def:kfeas}
Given diagnostics $\kappa_0^{\mathrm{emp}}(j)$ and frame bounds $(\alpha_j,\beta_j)$, we define the predicted feasibility bound
\begin{equation}\label{eq:kfeas-def}
K^{\mathrm{feas}}_j
\;:=\;
\kappa_0^{\mathrm{emp}}(j)\cdot \mathrm{cond}(\Gamma_j)\cdot m_j
\;=\;
\kappa_0^{\mathrm{emp}}(j)\cdot (\beta_j/\alpha_j)\cdot m_j.
\end{equation}
The main static theorem (T-SPT-00) justifies that $\kappa_j(v)$ is uniformly controlled on $v\in\mathcal{F}_j$ under the corresponding hypotheses, with explicit constants comparable to \eqref{eq:kfeas-def}.
\end{definition}

\subsection{Defect measures for robust invariance}

Exact feasibility invariance is often too strong, so we also quantify route mismatch between evolution and packaging.

\begin{definition}[Cellwise feasibility defect]\label{def:defect-infty}
For $u_j\in H_j$ define the cellwise defect ratio
\begin{equation}\label{eq:defect-infty}
\mathrm{def}^{\infty}_j(u_j)
\;:=\;
\max_{c\in\mathcal{C}_j}
\frac{\|P_{j,c}(I-E_j)u_j\|^2}{\|u_j\|^2},
\end{equation}
with $\mathrm{def}^{\infty}_j(0):=0$.
The robust persistence theorem (T-SPT-02) assumes a bound of the form
$\mathrm{def}^{\infty}_j(u_j(t)) \le \varepsilon_\infty f_j$ uniformly in time and depth, which is strictly stronger (and more relevant to localization) than a global bound on $\|(I-E_j)u_j(t)\|/\|u_j(t)\|$.
\end{definition}

\paragraph{Remark (finite-dimensional implementation).}
In the numerical experiments, each $H_j$ is represented as $\mathbb{R}^{n_j}$ with a contiguous block structure for cells and a synthesis matrix $\Psi_j\in\mathbb{R}^{n_j\times m_j}$ whose columns are the channels $\psi_{j,r}$.
The packaging projector $E_j$ is implemented via an orthonormal basis of $\Ran(\Psi_j)$, and all diagnostics in Definitions~\ref{def:kappa0-emp}--\ref{def:kfeas} are computed directly from $\Psi_j$.
 \section{Hypothesis Ledger (HL) for SPT-legality persistence}\label{sec:hypotheses}

This section records all assumptions as an explicit \emph{hypothesis ledger} (HL) in the style of \cite{Tsiokos2026SixBirds}.
The main theorems in Section~\ref{sec:main-theorems} will reference these slots verbatim.
The purpose is twofold:
(i) to make the closure regime explicit and mechanizable, and
(ii) to delineate sharp failure modes when a slot is dropped.

\subsection{Conventions for constants}

Throughout, $j\in\{0,\dots,J_{\max}\}$ denotes depth, $\mathcal{C}_j$ is the cell index set with baseline fraction $f_j=1/|\mathcal{C}_j|$, and $\Gamma_j:\mathbb{R}^{m_j}\to H_j$ is a channel map with feasible set $\mathcal{F}_j=\Ran(\Gamma_j)$ and packaging projector $E_j$ (Definition~\ref{def:packaging}).
For convenience we also define the global packaging projector $E:H\to H$ by
$Eu := \bigoplus_{j=0}^{J_{\max}} E_j(P_j u)$, whose range is $\bigoplus_j \mathcal{F}_j$.
We use the following artifact diagnostics (Definition~\ref{def:kappa0-emp}--\ref{def:kfeas}):
\[
\kappa^{\mathrm{emp}}_0(j),\qquad
\alpha_j,\beta_j,\qquad
\mathrm{cond}(\Gamma_j)=\beta_j/\alpha_j,\qquad
K^{\mathrm{feas}}_j := \kappa^{\mathrm{emp}}_0(j)\,\mathrm{cond}(\Gamma_j)\,m_j.
\]
We will also use a global envelope
\begin{equation}\label{eq:Kpred-global}
K^{\mathrm{pred}}_{\mathrm{glob}}
\;:=\;
\max_{0\le j\le J_{\max}} K^{\mathrm{feas}}_j.
\end{equation}

\subsection{Artifact slots (representation constraints)}

\paragraph{\textbf{HL-SPT-DELOC} (cell-delocalized channels).}
For each depth $j$, every channel $\psi_{j,r}$ is cell-delocalized with constant $\kappa_0(j)$ in the sense of \eqref{eq:deloc-hyp}, and we assume a uniform bound
\begin{equation}\label{eq:HL-DELOC}
\kappa^{\mathrm{emp}}_0(j)\le K_{\mathrm{deloc}}
\qquad \text{for all } j,
\end{equation}
for some finite $K_{\mathrm{deloc}}\ge 1$.
\begin{itemize}
\item \emph{Intended implication:} no single channel is wavepacket-like at the cell scale.
\item \emph{Failure modes if dropped:} feasibility can contain a cell-supported direction (FM-DELOC), making $\kappa_j$ achieve its maximum $1/f_j$ even statically.
\end{itemize}

\paragraph{\textbf{HL-SPT-FRAME} (nondegenerate frame bounds / conditioning).}
For each depth $j$, $\Gamma_j$ has positive frame lower bound $\alpha_j>0$ and finite condition proxy
\begin{equation}\label{eq:HL-FRAME}
\mathrm{cond}(\Gamma_j)=\beta_j/\alpha_j \le K_{\mathrm{cond}}
\qquad \text{for all } j,
\end{equation}
Equivalently, this slot requires stable synthesis (injective $\Gamma_j$); we intentionally separate it from the redundant/high-mode regime treated as a failure mode in FM-MODE.
for some finite $K_{\mathrm{cond}}\ge 1$.
\begin{itemize}
\item \emph{Intended implication:} coefficient energy is comparable to feasible state energy and near-dependence cannot amplify concentration in special coefficient directions.
\item \emph{Failure modes if dropped:} columns may be delocalized but extreme $\kappa_j(\Gamma_j a)$ events appear in the near-null directions of $\Gamma_j^\top\Gamma_j$ (FM-FRAME).
\end{itemize}

\paragraph{\textbf{HL-SPT-MODE} (controlled channel count / staged interface growth).}
The channel counts $m_j$ are controlled across depth, e.g.\ by a uniform mode envelope
\begin{equation}\label{eq:HL-MODE}
m_j \le m_{\max}\qquad\text{for all } j,
\end{equation}
or more generally by a specified growth law $m_j\le M(j)$ with $M(j)$ at most polynomial in $j$.
\begin{itemize}
\item \emph{Intended implication:} interface complexity does not explode with depth, enabling depth-uniform legality thresholds.
\item \emph{Failure modes if dropped:} when $m_j$ grows toward the full band dimension, feasibility approximates all of $H_j$ and a fixed $K_\kappa$ cannot hold uniformly (FM-MODE).
\end{itemize}

\subsection{System slots (compatibility of dynamics with feasibility)}

\paragraph{\textbf{HL-SPT-INV} (exact feasibility invariance).}
The dynamical system is \emph{SPT-compatible} in the strong sense that feasibility is invariant bandwise:
\begin{equation}\label{eq:HL-INV}
u_j(0)\in \mathcal{F}_j \;\Rightarrow\; u_j(t)\in \mathcal{F}_j \quad \text{for all } t \text{ in the lifespan and all } j.
\end{equation}
Equivalently, if $u_j(0)$ is packaged then $E_j u_j(t)=u_j(t)$ for all later times.
\begin{itemize}
\item \emph{Intended implication:} the evolution respects the packaged representation (P5/P2 compatibility).
\item \emph{Failure modes if dropped:} even tiny leakage out of feasibility can cause deep-band $\kappa_j$ spikes (FM-INV).
\end{itemize}

\begin{remark}[HL-SPT-INV as an invariance/commutation contract]\label{rem:inv-contract}
Let $\Phi_t:H\to H$ denote the evolution map (when it exists) so that $u(t)=\Phi_t(u(0))$.
HL-SPT-INV is equivalently the invariance of the packaged subspace $\Ran(E)$:
\[
u(0)=Eu(0)\ \Rightarrow\ u(t)=Eu(t)\quad\text{for all }t,
\]
or, in operator form,
\[
(I-E)\,\Phi_t\,E = 0 \qquad \text{for all } t \text{ in the lifespan.}
\]
In our toy ladder model (Section~\ref{sec:toy-models}) this holds by construction; in general systems it is a nontrivial compatibility slot rather than a generic consequence of accounting.
\end{remark}

\paragraph{\textbf{HL-SPT-ROBUST$\infty$} (robust feasibility with cellwise defect control).}
A weaker but more realistic hypothesis allows leakage, provided it is controlled at the cell scale:
there exists $\varepsilon_\infty\ge 0$ such that for all $j$ and all times $t$,
\begin{equation}\label{eq:HL-ROBUST}
\mathrm{def}^{\infty}_j(u_j(t))
\;=\;
\max_{c\in\mathcal{C}_j}\frac{\|P_{j,c}(I-E_j)u_j(t)\|^2}{\|u_j(t)\|^2}
\;\le\;
\varepsilon_\infty\, f_j.
\end{equation}
\begin{itemize}
\item \emph{Intended implication:} route mismatch between evolution and packaging is locally small in the same partition used to define legality.
\item \emph{Failure modes if dropped:} global leakage $\|(I-E_j)u_j\|/\|u_j\|$ may be tiny while deep-band $\kappa_j$ approaches $1/f_j$ (NG-DIST).
\end{itemize}

\begin{remark}[HL-SPT-ROBUST$\infty$ as a cellwise route-mismatch bound]\label{rem:robust-contract}
If $\Phi_t$ denotes the evolution map, HL-SPT-ROBUST$\infty$ can be viewed as a \emph{cellwise} route-mismatch contract:
the off-packaged component $(I-E_j)u_j(t)$ is allowed, but its per-cell energy fraction relative to $\|u_j(t)\|^2$ is bounded by $\varepsilon_\infty f_j$.
This relative (ratio) formulation is intentional: deep bands may carry tiny global norm, so absolute leakage bounds do not control the max-cell audit (NG-DIST).
\end{remark}

\subsection{Audit hook and interpretation}

The HL slots are designed to be checkable.
In the accompanying codebase, every experiment run produces an \texttt{audit.json} file recording:
(i) whether HL-SPT-DELOC and HL-SPT-FRAME pass numerical thresholds (from $\kappa^{\mathrm{emp}}_0$ and conditioning),
(ii) the realized mode growth profile $m_j$ (HL-SPT-MODE), and
(iii) evidence for invariance or robust defect control when applicable (HL-SPT-INV / HL-SPT-ROBUST$\infty$).
The ``FM'' and ``NG'' experiments intentionally violate specific slots; the audit report is part of the regression harness that prevents silent drift while the paper is written.
 \section{Main Theorems (statements)}\label{sec:main-theorems}

We now state the persistence results in ``Six-Birds theorem'' form: each result explicitly references the hypothesis ledger (HL) slots from Section~\ref{sec:hypotheses}, specifies the constants, and lists sharp failure modes when a slot is dropped.
Proofs are given in Sections~6--8.

\subsection{T-SPT-00: feasibility geometry implies legality bounds}

\begin{theorem}[T-SPT-00: feasibility legality bound from channel delocalization]\label{thm:t-spt-00}
Fix a depth $j\le J_{\max}$ and a channel map $\Gamma_j:\mathbb{R}^{m_j}\to H_j$ with feasible set $\mathcal{F}_j=\Ran(\Gamma_j)$ and cell partition $\mathcal{C}_j$.
Assume the artifact slots:
\begin{itemize}
\item \textbf{HL-SPT-DELOC} at depth $j$, with empirical bound $\kappa^{\mathrm{emp}}_0(j)\le K_{\mathrm{deloc}}$,
\item \textbf{HL-SPT-FRAME} at depth $j$, with frame bounds $0<\alpha_j\le \beta_j$ and $\mathrm{cond}(\Gamma_j)=\beta_j/\alpha_j\le K_{\mathrm{cond}}$.
\end{itemize}
Then every nonzero feasible vector $v\in\mathcal{F}_j$ satisfies
\begin{equation}\label{eq:t-spt-00-bound}
\kappa_j(v)
\;\le\;
K^{\mathrm{feas}}_j
\;:=\;
\kappa^{\mathrm{emp}}_0(j)\cdot \mathrm{cond}(\Gamma_j)\cdot m_j
\;\le\;
K_{\mathrm{deloc}}\,K_{\mathrm{cond}}\,m_j.
\end{equation}

\paragraph{Failure modes.}
If HL-SPT-DELOC is dropped, $\mathcal{F}_j$ may contain cell-supported directions and $\kappa_j$ can reach its maximum $1/f_j$ even statically (FM-DELOC).
If HL-SPT-FRAME is dropped, channels may be individually delocalized while near-dependence yields coefficient directions $a$ with extreme $\kappa_j(\Gamma_j a)$ tails (FM-FRAME).
\end{theorem}

\subsection{T-SPT-01: exact invariance implies legality persistence}

\begin{theorem}[T-SPT-01: legality persistence from exact feasibility invariance]\label{thm:t-spt-01}
Assume the artifact slots \textbf{HL-SPT-DELOC}, \textbf{HL-SPT-FRAME}, and \textbf{HL-SPT-MODE}, so that $K^{\mathrm{pred}}_{\mathrm{glob}}$ in \eqref{eq:Kpred-global} is finite.
Assume also the system slot \textbf{HL-SPT-INV} (exact feasibility invariance).

Let $u(t)\in H$ be a trajectory on its lifespan such that $u_j(0)\in \mathcal{F}_j$ for all $j$.
Then for every time $t$ in the lifespan and every depth $j$ with $u_j(t)\neq 0$,
\begin{equation}\label{eq:t-spt-01-bound}
\kappa_j(u_j(t)) \;\le\; K^{\mathrm{feas}}_j
\qquad \text{and hence} \qquad
\max_{j}\kappa_j(u_j(t)) \;\le\; K^{\mathrm{pred}}_{\mathrm{glob}}.
\end{equation}
In particular, the trajectory remains SPT-legal for all time with legality threshold
\begin{equation}\label{eq:t-spt-01-K}
K_\kappa := K^{\mathrm{pred}}_{\mathrm{glob}}.
\end{equation}

\paragraph{Failure modes.}
If HL-SPT-INV is dropped, even arbitrarily small leakage out of feasibility can induce deep-band $\kappa_j$ spikes and immediate illegality (FM-INV).
If HL-SPT-MODE is dropped, $K^{\mathrm{pred}}_{\mathrm{glob}}$ may grow without bound with depth (FM-MODE), so a fixed legality threshold is not meaningful.

\paragraph{Interpretation (closure theorem).}
T-SPT-01 is intentionally conditional: it states that once the system-level invariance slot holds, the legal regime is forward-closed with explicit constants determined by the artifact.
The theorem does not claim that invariance holds for arbitrary dynamics; verifying HL-SPT-INV (or its robust substitute) is the principal system-side task.
\end{theorem}

\subsection{T-SPT-02: robust persistence under cellwise defect control}

\begin{theorem}[T-SPT-02: robust legality persistence from cellwise feasibility defect]\label{thm:t-spt-02}
Assume the artifact slots \textbf{HL-SPT-DELOC}, \textbf{HL-SPT-FRAME}, and \textbf{HL-SPT-MODE}, and let $K^{\mathrm{pred}}_{\mathrm{glob}}$ be as in \eqref{eq:Kpred-global}.
Assume the robust system slot \textbf{HL-SPT-ROBUST$\infty$}: there exists $\varepsilon_\infty\ge 0$ such that for all depths $j$ and all times $t$ in the lifespan,
\[
\mathrm{def}^{\infty}_j(u_j(t)) \;\le\; \varepsilon_\infty f_j.
\]
Then the trajectory is SPT-legal for all time with the explicit robust legality threshold
\begin{equation}\label{eq:t-spt-02-Krob}
K^{\mathrm{rob}}(\varepsilon_\infty)
\;:=\;
\Big(\sqrt{K^{\mathrm{pred}}_{\mathrm{glob}}} + \sqrt{\varepsilon_\infty}\Big)^2.
\end{equation}
More precisely, for every time $t$ and every depth $j$ with $u_j(t)\neq 0$,
\begin{equation}\label{eq:t-spt-02-bound}
\kappa_j(u_j(t)) \;\le\; K^{\mathrm{rob}}(\varepsilon_\infty).
\end{equation}

\paragraph{Failure modes.}
If HL-SPT-ROBUST$\infty$ is replaced by a purely global leakage bound (e.g.\ $\|(I-E_j)u_j\|/\|u_j\|\ll 1$), robust persistence can fail at deep scales because tiny global leakage can concentrate in a single cell (NG-DIST).
If HL-SPT-DELOC/FRAME/MODE are dropped, the feasibility geometry itself may admit near-maximal concentration or may produce extreme tails (FM-DELOC, FM-FRAME, FM-MODE), and no uniform robust legality threshold exists.

\paragraph{Interpretation (why the defect is formulated at cell resolution).}
Because legality is defined by a max-cell functional, robust persistence cannot be phrased using only a global distance-to-feasibility.
HL-SPT-ROBUST$\infty$ is designed to be the minimal natural robustness slot aligned with the audit, and NG-DIST shows that weakening it to a global leakage metric is insufficient.
\end{theorem}

\paragraph{Remark (interpretation).}
T-SPT-01 and T-SPT-02 express closure of the P6-defined legal regime (bounded $\kappa_j$) under P5/P2/P4 compatibility hypotheses in the sense of \cite{Tsiokos2026SixBirds}.
In the Six-Birds language, HL-SPT-INV or HL-SPT-ROBUST$\infty$ is the explicit closure slot (the analogue of a ``macro closure'' requirement), while the artifact slots enforce structural anti-localization rather than relying on accounting alone.
 \section{Proof of T-SPT-00}\label{sec:proof-t-spt-00}

We prove Theorem~\ref{thm:t-spt-00}.
Fix a depth $j\le J_{\max}$ and suppress the depth index $j$ in the notation to reduce clutter:
write $\mathcal{C}=\mathcal{C}_j$, $f=f_j$, $\Gamma=\Gamma_j$, $m=m_j$, and $\psi_r=\psi_{j,r}$.
Let $P_c:=P_{j,c}$ denote the cell projectors.

If $v=0$ then $\kappa_j(v)=0$ by convention (Definitions~\ref{def:cell-conc}--\ref{def:kappa-def}), so assume $v\neq 0$.
Since $v\in \mathcal{F}_j=\Ran(\Gamma)$, there exists $a\in\mathbb{R}^m$ such that
\[
v=\Gamma a=\sum_{r=1}^m a_r \psi_r.
\]

\paragraph{Step 1: bound the energy in a fixed cell.}
Fix a cell $c\in\mathcal{C}$.
By linearity,
\[
P_c v \;=\; \sum_{r=1}^m a_r\,P_c\psi_r.
\]
Using Cauchy--Schwarz in $\mathbb{R}^m$,
\begin{align}
\|P_c v\|
&\le \sum_{r=1}^m |a_r|\,\|P_c\psi_r\|
\le \Big(\sum_{r=1}^m |a_r|^2\Big)^{1/2}
     \Big(\sum_{r=1}^m \|P_c\psi_r\|^2\Big)^{1/2}
= \|a\|\Big(\sum_{r=1}^m \|P_c\psi_r\|^2\Big)^{1/2}.
\end{align}
Squaring yields
\begin{equation}\label{eq:proof-cell-bound-1}
\|P_c v\|^2 \;\le\; \|a\|^2 \sum_{r=1}^m \|P_c\psi_r\|^2.
\end{equation}

\paragraph{Step 2: use empirical delocalization to control the channel energy per cell.}
By Definition~\ref{def:kappa0-emp} (empirical delocalization constant), for every $r$ and every cell $c$,
\[
\|P_c\psi_r\|^2 \;\le\; \kappa_0^{\mathrm{emp}}(j)\,f\,\|\psi_r\|^2.
\]
Summing over $r$ gives
\begin{equation}\label{eq:proof-cell-bound-2}
\sum_{r=1}^m \|P_c\psi_r\|^2
\;\le\;
\kappa_0^{\mathrm{emp}}(j)\,f\sum_{r=1}^m \|\psi_r\|^2.
\end{equation}

\paragraph{Step 3: relate $\sum_r\|\psi_r\|^2$ to the frame upper bound.}
Let $G:=\Gamma^\top\Gamma\in\mathbb{R}^{m\times m}$ be the Gram matrix.
Then $\mathrm{tr}(G)=\sum_{r=1}^m \|\psi_r\|^2$.
Under HL-SPT-FRAME at depth $j$ (Definition~\ref{def:frame-bounds}), the eigenvalues of $G$ lie in $[\alpha_j,\beta_j]$, hence
\begin{equation}\label{eq:proof-trace-bound}
\sum_{r=1}^m \|\psi_r\|^2 \;=\; \mathrm{tr}(G) \;\le\; m\,\beta_j.
\end{equation}

\paragraph{Step 4: combine and normalize by $\|v\|^2$.}
Combining \eqref{eq:proof-cell-bound-1}, \eqref{eq:proof-cell-bound-2}, and \eqref{eq:proof-trace-bound}, we obtain
\[
\|P_c v\|^2
\;\le\;
\|a\|^2\cdot \kappa_0^{\mathrm{emp}}(j)\,f\cdot m\,\beta_j.
\]
On the other hand, by the frame lower bound in \eqref{eq:frame-bounds},
\[
\|v\|^2=\|\Gamma a\|^2 \;\ge\; \alpha_j\,\|a\|^2.
\]
Therefore the cell energy fraction satisfies
\[
\frac{\|P_c v\|^2}{\|v\|^2}
\;\le\;
\kappa_0^{\mathrm{emp}}(j)\,f\,m\,\frac{\beta_j}{\alpha_j}
\;=\;
\kappa_0^{\mathrm{emp}}(j)\,f\,m\,\mathrm{cond}(\Gamma_j).
\]
Taking the maximum over $c\in\mathcal{C}$ yields (Definition~\ref{def:cell-conc})
\[
\conc_j(v) \;\le\; \kappa_0^{\mathrm{emp}}(j)\,f_j\,m_j\,\mathrm{cond}(\Gamma_j).
\]
Dividing by $f_j$ gives (Definition~\ref{def:kappa-def})
\[
\kappa_j(v)
\;=\;
\frac{\conc_j(v)}{f_j}
\;\le\;
\kappa_0^{\mathrm{emp}}(j)\,\mathrm{cond}(\Gamma_j)\,m_j
\;=\;
K^{\mathrm{feas}}_j,
\]
which is exactly \eqref{eq:t-spt-00-bound} and completes the proof.

\paragraph{Remark (where $m_j$ enters).}
The factor $m_j$ arises from the trace bound \eqref{eq:proof-trace-bound}.
If the channels are orthonormal (or a tight frame), one may replace this step with a sharper identity, but the present proof isolates the two genuinely necessary artifact controls:
cell-delocalization of columns (HL-SPT-DELOC) and nondegenerate frame bounds (HL-SPT-FRAME).
 \section{Proof of T-SPT-01}\label{sec:proof-t-spt-01}

We prove Theorem~\ref{thm:t-spt-01}.
Fix a trajectory $u(t)\in H$ on its lifespan and assume $u_j(0)\in\mathcal{F}_j$ for all depths $j$.
By the system slot HL-SPT-INV (exact feasibility invariance, \eqref{eq:HL-INV}), it follows that
\[
u_j(t)\in \mathcal{F}_j \qquad \text{for all times $t$ in the lifespan and all depths $j$.}
\]
For each fixed $t$ and each depth $j$ with $u_j(t)\neq 0$, we may therefore apply the static feasibility bound (Theorem~\ref{thm:t-spt-00}) to the feasible vector $v=u_j(t)\in\mathcal{F}_j$ to obtain
\[
\kappa_j(u_j(t)) \;\le\; K^{\mathrm{feas}}_j.
\]
Taking the maximum over depths yields
\[
\max_{0\le j\le J_{\max}}\kappa_j(u_j(t))
\;\le\;
\max_{0\le j\le J_{\max}} K^{\mathrm{feas}}_j
\;=\;
K^{\mathrm{pred}}_{\mathrm{glob}},
\]
where the last equality is the definition \eqref{eq:Kpred-global}.
This is exactly \eqref{eq:t-spt-01-bound} and implies SPT-legality with threshold $K_\kappa=K^{\mathrm{pred}}_{\mathrm{glob}}$ as stated in \eqref{eq:t-spt-01-K}.
 \section{Proof of T-SPT-02}\label{sec:proof-t-spt-02}

We prove Theorem~\ref{thm:t-spt-02}.
Fix a time $t$ in the lifespan and a depth $j\le J_{\max}$.
Write $u:=u_j(t)\in H_j$ and decompose it into feasible and off-feasible parts using the packaging projector $E_j$ (Definition~\ref{def:packaging}):
\[
u \;=\; E_j u \;+\; (I-E_j)u \;=:\; v \;+\; w,
\qquad v\in\mathcal{F}_j,\quad w\perp \mathcal{F}_j.
\]
If $u=0$ then $\kappa_j(u)=0$ by convention, so assume $u\neq 0$.

\paragraph{Step 1: cellwise control of the off-feasible component.}
By the robust system slot HL-SPT-ROBUST$\infty$ \eqref{eq:HL-ROBUST} and the definition of $\mathrm{def}^\infty_j$ (Definition~\ref{def:defect-infty}),
for every cell $c\in\mathcal{C}_j$ we have
\begin{equation}\label{eq:proof-robust-cellwise}
\|P_{j,c}w\|^2
\;=\;
\|P_{j,c}(I-E_j)u\|^2
\;\le\;
\varepsilon_\infty f_j \,\|u\|^2.
\end{equation}
In particular,
\begin{equation}\label{eq:proof-robust-cellwise-sqrt}
\|P_{j,c}w\|
\;\le\;
\sqrt{\varepsilon_\infty f_j}\,\|u\|
\qquad \text{for all } c.
\end{equation}

\paragraph{Step 2: control the feasible component using T-SPT-00.}
Since $v=E_j u\in\mathcal{F}_j$, we may apply T-SPT-00 (Theorem~\ref{thm:t-spt-00}) to obtain
\[
\kappa_j(v)\le K^{\mathrm{feas}}_j \le K^{\mathrm{pred}}_{\mathrm{glob}},
\]
where the last inequality uses the definition \eqref{eq:Kpred-global}.
Unpacking the definition of $\kappa_j$ (Definition~\ref{def:kappa-def}), this means that for every cell $c$,
\begin{equation}\label{eq:proof-feasible-cell}
\|P_{j,c}v\|^2 \;\le\; f_j\,K^{\mathrm{pred}}_{\mathrm{glob}}\,\|v\|^2
\qquad\Rightarrow\qquad
\|P_{j,c}v\|\;\le\;\sqrt{f_j\,K^{\mathrm{pred}}_{\mathrm{glob}}}\,\|v\|.
\end{equation}
Finally note that $\|v\|=\|E_j u\|\le \|u\|$ since $E_j$ is an orthogonal projector.
Hence
\begin{equation}\label{eq:proof-feasible-cell-2}
\|P_{j,c}v\|
\;\le\;
\sqrt{f_j\,K^{\mathrm{pred}}_{\mathrm{glob}}}\,\|u\|
\qquad \text{for all } c.
\end{equation}

\paragraph{Step 3: combine feasible and off-feasible parts.}
By the triangle inequality in $H_{j,c}$,
\[
\|P_{j,c}u\|
\;=\;
\|P_{j,c}(v+w)\|
\;\le\;
\|P_{j,c}v\|+\|P_{j,c}w\|.
\]
Using \eqref{eq:proof-robust-cellwise-sqrt} and \eqref{eq:proof-feasible-cell-2}, we get
\[
\|P_{j,c}u\|
\;\le\;
\Big(\sqrt{f_j\,K^{\mathrm{pred}}_{\mathrm{glob}}} + \sqrt{\varepsilon_\infty f_j}\Big)\,\|u\|
\;=\;
\sqrt{f_j}\Big(\sqrt{K^{\mathrm{pred}}_{\mathrm{glob}}} + \sqrt{\varepsilon_\infty}\Big)\,\|u\|.
\]
Squaring and dividing by $\|u\|^2$ yields, for every cell $c$,
\[
\frac{\|P_{j,c}u\|^2}{\|u\|^2}
\;\le\;
f_j\Big(\sqrt{K^{\mathrm{pred}}_{\mathrm{glob}}} + \sqrt{\varepsilon_\infty}\Big)^2.
\]
Taking the maximum over $c$ gives
\[
\conc_j(u)\;\le\; f_j\Big(\sqrt{K^{\mathrm{pred}}_{\mathrm{glob}}} + \sqrt{\varepsilon_\infty}\Big)^2.
\]
Dividing by $f_j$ (Definition~\ref{def:kappa-def}) yields
\[
\kappa_j(u)
\;\le\;
\Big(\sqrt{K^{\mathrm{pred}}_{\mathrm{glob}}} + \sqrt{\varepsilon_\infty}\Big)^2
\;=\;
K^{\mathrm{rob}}(\varepsilon_\infty),
\]
which is exactly \eqref{eq:t-spt-02-bound}--\eqref{eq:t-spt-02-Krob}.
Since $t$ and $j$ were arbitrary, the bound holds for all times and depths, proving T-SPT-02.

\paragraph{Remark (why defect must be cellwise).}
The key step is the uniform per-cell bound \eqref{eq:proof-robust-cellwise}.
A purely global leakage control such as $\|w\|/\|u\|\ll 1$ does \emph{not} imply any nontrivial bound on $\max_c\|P_{j,c}w\|^2/\|u\|^2$ when $f_j$ is small (deep refinement).
Indeed, it is possible for $\|w\|/\|u\|$ to be tiny while $w$ is supported in a single cell, in which case it can dominate $\conc_j(u)$ and drive $\kappa_j(u)$ close to its maximum $1/f_j$; this is the phenomenon isolated in NG-DIST.
 \section{Remarks and Non-implications}\label{sec:nonimplications}

This section records several statements that may sound plausible on first reading but are \emph{not} implied by our definitions and theorems.
The point is not pedantry: each non-implication corresponds to a concrete failure mode demonstrated by a minimal counterexample and a numerical experiment in Section~12.
In the Six-Birds framing \cite{Tsiokos2026SixBirds}, these are exactly the points where one is tempted to ``smuggle in'' closure using accounting (P6) when what is needed is structural compatibility (P2/P5/P4) or an explicit defect control (P3).

\subsection{Ledger invariants do not imply legality}

\paragraph{Non-implication 1 (energy does not control $\kappa$).}
Even if a trajectory preserves a global ``ledger'' quantity such as energy $\|u(t)\|^2$, this does not prevent the cell concentration audit $\kappa_j(u_j(t))$ (Definition~\ref{def:kappa-def}) from becoming large.
The obstruction is simple: if the feasible set $\mathcal{F}_j$ contains a localized direction, then energy-preserving motion \emph{within} $\mathcal{F}_j$ can rotate mass from a delocalized feasible state to a localized feasible state, driving $\kappa_j$ close to its ceiling $1/f_j$ while $\|u\|^2$ remains constant.

\paragraph{What would be needed instead.}
To prevent this, one needs a structural anti-localization mechanism: either (i) feasibility itself excludes localized directions (HL-SPT-DELOC), or (ii) legality is defined by a different audit that is actually controlled by the ledger.
Our no-go construction NG-LEDGER (Section~12) makes this explicit: $\|u(t)\|^2$ is flat to numerical precision while $\kappa_j(u_j(t))$ reaches $1/f_j$.

\subsection{Delocalized columns do not protect combinations without frame control}

\paragraph{Non-implication 2 (column delocalization $\not\Rightarrow$ combination delocalization).}
A common misconception is that if every channel $\psi_{j,r}$ is individually cell-delocalized, then every feasible superposition $v=\Gamma_j a$ must also be cell-delocalized.
This is false without a quantitative frame lower bound.
Theorem~\ref{thm:t-spt-00} makes the dependence explicit: the bound scales with $\mathrm{cond}(\Gamma_j)=\beta_j/\alpha_j$ (Definition~\ref{def:frame-bounds}).
When $\alpha_j$ is small (near-dependence of channels), there exist coefficient directions $a$ that produce superpositions with large localized peaks, despite each column being spread out.

\paragraph{What would be needed instead.}
One must control both HL-SPT-DELOC and HL-SPT-FRAME.
The FM-FRAME experiment (Section~12) isolates this phenomenon: $\kappa_0^{\mathrm{emp}}$ remains bounded (columns are delocalized) while the conditioning worsens and the tail events of $\kappa_j(\Gamma_j a)$ explode.

\subsection{Global distance to feasibility is not a robustness notion for $\kappa$}

\paragraph{Non-implication 3 (small global leakage $\not\Rightarrow$ small $\kappa$).}
In robust settings, it is natural to measure feasibility mismatch by a global ratio
$\|(I-E_j)u_j\|/\|u_j\|$.
However, this is not a robustness notion for a \emph{max-cell} audit.
A small global leakage can still concentrate into a single cell and dominate $\mathrm{conc}_j(u_j)$, especially at deep scales where $f_j=1/|\mathcal{C}_j|$ is small.
The robust persistence theorem (T-SPT-02) therefore assumes the stronger \emph{cellwise} defect control
$\mathrm{def}^{\infty}_j(u_j(t))\le \varepsilon_\infty f_j$ (Definition~\ref{def:defect-infty} and HL-SPT-ROBUST$\infty$ \eqref{eq:HL-ROBUST}).

\paragraph{What would be needed instead.}
If one insists on a global leakage metric, one must supplement it with additional structure that links global leakage to cellwise leakage (for instance, an \emph{a priori} delocalization bound on the off-feasible component).
Absent such a link, robust legality can fail even when the global leakage is arbitrarily small.
The NG-DIST experiment (Section~12) demonstrates precisely this gap: leakage ratios below $10^{-3}$ coexist with deep-band $\kappa_j$ near the maximal value $1/f_j$.

\subsection{Summary}

The previous points identify the genuine ``hard slots'' behind legality persistence:
anti-localization must be enforced structurally in feasibility (HL-SPT-DELOC), feasible superpositions must be stable under coefficients (HL-SPT-FRAME), depth-uniform thresholds require controlled interface growth (HL-SPT-MODE), and robustness requires defect bounds formulated at the same resolution as the audit itself (HL-SPT-ROBUST$\infty$).
These observations motivate the sharp no-go constructions in Section~10 and the experiment suite in Section~12.
 \section{No-go results and sharp counterexamples}\label{sec:nogo}

This section sharpens the boundary of legality persistence by giving minimal counterexamples.
Each counterexample corresponds to dropping (or weakening) one HL slot from Section~\ref{sec:hypotheses}.
These are not PDE pathologies: all constructions are finite-dimensional, expressed directly in the Hilbert/band/\allowbreak cell model of Section~\ref{sec:definitions}.
Each item also has a numerical instantiation in Section~12, with a frozen figure in \texttt{paper/figures/}.

\paragraph{Roadmap.}
FM-DELOC shows why HL-SPT-DELOC is necessary (Prop.~\ref{prop:fm-deloc});
FM-FRAME shows why HL-SPT-FRAME is necessary even when columns are delocalized (Prop.~\ref{prop:fm-frame});
FM-MODE shows why depth-uniform thresholds require controlled interface growth (Prop.~\ref{prop:fm-mode});
NG-LEDGER shows why ledger invariants do not prevent localization blow-up (Prop.~\ref{prop:ng-ledger});
NG-DIST shows why global distance-to-feasibility is not a robustness notion for $\kappa$ (Prop.~\ref{prop:ng-dist});
and FM-INV packages NG-DIST into a dynamical ``leakage breaks legality'' mechanism (Prop.~\ref{prop:fm-inv}).

\subsection{FM-DELOC: feasibility contains wavepackets}

\begin{proposition}[FM-DELOC: if feasibility contains a cell-supported direction, legality fails statically]\label{prop:fm-deloc}
Fix a depth $j$ and cell partition $\mathcal{C}_j$ with baseline fraction $f_j$.
Assume the feasible set $\mathcal{F}_j\subseteq H_j$ contains a nonzero vector $v$ supported in a single cell, i.e.\ there exists $c_\star\in\mathcal{C}_j$ such that
\[
P_{j,c_\star}v=v
\qquad \text{and hence } P_{j,c}v=0 \text{ for } c\neq c_\star.
\]
Then $\conc_j(v)=1$ and therefore
\[
\kappa_j(v)=\frac{1}{f_j}.
\]
In particular, for any threshold $K_\kappa<0.9/f_j$, the feasible set contains a vector violating SPT-legality by a factor at least $0.9/f_j$ (indeed, by the full factor $1/f_j$).
\end{proposition}

\paragraph{Proof sketch.}
All energy lies in one cell, so the maximizing cell fraction in Definition~\ref{def:cell-conc} equals $1$; dividing by $f_j$ gives $\kappa_j(v)=1/f_j$ (Definition~\ref{def:kappa-def}).

\paragraph{HL interpretation.}
This is exactly what HL-SPT-DELOC rules out: if even one channel is cell-supported, then the feasible set can contain wavepacket-like directions.

\paragraph{Experiment pointer.}
FM-DELOC is instantiated by a feasible channel that is literally cell-supported; see \nolinkurl{paper/figures/localized_channel_cell_energy.png}.

\subsection{FM-FRAME: delocalized columns without frame control do not protect combinations}

\begin{proposition}[FM-FRAME: bounded per-column delocalization does not prevent extreme concentration without a frame lower bound]\label{prop:fm-frame}
Fix a depth $j$ and let $f_j=1/|\mathcal{C}_j|$.
There exist channel maps $\Gamma_j:\mathbb{R}^{m_j}\to H_j$ with the following properties simultaneously:
\begin{enumerate}
\item (\emph{Columns delocalized}) $\kappa^{\mathrm{emp}}_0(j)\le 2$ (Definition~\ref{def:kappa0-emp}).
\item (\emph{Ill-conditioned}) $\mathrm{cond}(\Gamma_j)$ is arbitrarily large (Definition~\ref{def:frame-bounds}).
\item (\emph{Feasible concentration}) the feasible set $\mathcal{F}_j=\Ran(\Gamma_j)$ contains a nonzero $v$ with
\[
\kappa_j(v)=\frac{1}{f_j}\;\ge\; 0.9/f_j.
\]
\end{enumerate}
(See \cite{Christensen2016Frames} for standard conditioning/frame stability notions.)
\end{proposition}

\paragraph{Proof sketch (explicit two-channel construction).}
Let $m_j=2$ and choose orthonormal vectors $u,e\in H_j$ such that:
(i) $u$ is perfectly cell-uniform (so $\kappa_j(u)=1$), and
(ii) $e$ is supported in a single cell $c_\star$ (so $\kappa_j(e)=1/f_j$),
with the additional choice that $P_{j,c_\star}u\perp e$ (e.g.\ take $e$ in a different internal coordinate within the same cell when $\dim(H_{j,c_\star})\ge 2$).
For $\varepsilon>0$, define channels
\[
\psi_1 := u + \varepsilon e,
\qquad
\psi_2 := u - \varepsilon e,
\qquad
\Gamma_j a := a_1\psi_1+a_2\psi_2.
\]
Because $P_{j,c_\star}u\perp e$, we have $\|P_{j,c_\star}\psi_r\|^2=\|P_{j,c_\star}u\|^2+\varepsilon^2\|e\|^2=f_j\|u\|^2+\varepsilon^2$ and $\|\psi_r\|^2=\|u\|^2+\varepsilon^2\|e\|^2$.
If $\|u\|=\|e\|=1$ and $\varepsilon^2\le f_j$, then
\[
\max_c\frac{\|P_{j,c}\psi_r\|^2}{f_j\|\psi_r\|^2}
=
\frac{f_j+\varepsilon^2}{f_j(1+\varepsilon^2)}
\le 2,
\]
so $\kappa_0^{\mathrm{emp}}(j)\le 2$.
However, the Gram matrix has eigenvalues $2$ and $2\varepsilon^2$, hence $\mathrm{cond}(\Gamma_j)=1/\varepsilon^2$ is arbitrarily large.
Finally, choosing coefficients $a=(1,-1)$ gives $\Gamma_j a = 2\varepsilon e$, which is cell-supported and therefore satisfies $\kappa_j(\Gamma_j a)=1/f_j$.
This shows that delocalized columns alone do not protect feasible superpositions from localization; one needs HL-SPT-FRAME.

\paragraph{Experiment pointer.}
FM-FRAME is instantiated by constructing delocalized columns with increasingly ill-conditioned Gram matrices and measuring concentration tails; see \nolinkurl{paper/figures/kappa_tail_vs_cond.png}.

\subsection{FM-MODE: mode explosion destroys depth-uniform legality}

\begin{proposition}[FM-MODE: without mode-growth control, no depth-uniform legality threshold exists]\label{prop:fm-mode}
Consider a family of depths with $|\mathcal{C}_j|\to\infty$ (hence $f_j\to 0$).
Suppose there exist depths $j$ for which the feasible set is all of the band, $\mathcal{F}_j=H_j$ (e.g.\ $\Gamma_j$ is surjective, which is possible when $m_j\ge \dim(H_j)$).
Then for each such $j$ there exists a feasible vector $v\in\mathcal{F}_j$ with
\[
\kappa_j(v)=\frac{1}{f_j}.
\]
Consequently, any fixed legality threshold $K_\kappa$ independent of $j$ fails for sufficiently deep $j$.

More generally, even in the best-case artifact regime where $\kappa^{\mathrm{emp}}_0(j)$ and $\mathrm{cond}(\Gamma_j)$ remain bounded, the feasibility envelope in T-SPT-00 scales linearly in $m_j$:
\[
\sup_{v\in\mathcal{F}_j\setminus\{0\}}\kappa_j(v)
\;\lesssim\;
\kappa^{\mathrm{emp}}_0(j)\,\mathrm{cond}(\Gamma_j)\,m_j,
\]
so if $m_j$ grows without bound with depth then the best available legality bound also grows without bound.
\end{proposition}

\paragraph{Proof sketch.}
If $\mathcal{F}_j=H_j$, then $H_j$ contains a cell-supported vector (choose any nonzero element of $H_{j,c_\star}$), which achieves $\kappa_j=1/f_j$ by the same argument as in Prop.~\ref{prop:fm-deloc}.
The final inequality is just Theorem~\ref{thm:t-spt-00} written as a uniform bound on feasible vectors.

\paragraph{Experiment pointer.}
FM-MODE is instantiated by running the invariant ladder dynamics under different growth laws $m_j$; see \nolinkurl{paper/figures/max_kappa_vs_depth_by_regime.png}.

\subsection{NG-LEDGER: energy conservation does not prevent concentration blow-up}

\begin{proposition}[NG-LEDGER: an energy-preserving evolution can drive $\kappa_j$ to the ceiling]\label{prop:ng-ledger}
Fix a depth $j$.
Assume the feasible set $\mathcal{F}_j$ contains two orthonormal vectors $\psi_{\mathrm{uni}},\psi_{\mathrm{loc}}\in \mathcal{F}_j$ such that
\[
\kappa_j(\psi_{\mathrm{uni}})=1,
\qquad
\kappa_j(\psi_{\mathrm{loc}})=\frac{1}{f_j}.
\]
Define an energy-preserving evolution in the two-dimensional subspace $\mathrm{span}\{\psi_{\mathrm{uni}},\psi_{\mathrm{loc}}\}$ by
\[
u(\theta)=\cos\theta\,\psi_{\mathrm{uni}}+\sin\theta\,\psi_{\mathrm{loc}}.
\]
Then $\|u(\theta)\|^2\equiv 1$ for all $\theta$, but $\kappa_j(u(\theta))$ can be made arbitrarily close to $1/f_j$.
In particular, there exists $\theta$ such that
\[
\kappa_j(u(\theta)) \ge 0.9/f_j.
\]
\end{proposition}

\paragraph{Proof sketch.}
Energy is constant because the evolution is a rotation in an orthonormal basis.
As $\theta\to \pi/2$, $u(\theta)\to \psi_{\mathrm{loc}}$ in norm, and $\kappa_j(\cdot)$ is continuous on $H_j\setminus\{0\}$ because it is the maximum of finitely many continuous functions
$v\mapsto \|P_{j,c}v\|^2/\|v\|^2$.
Thus $\kappa_j(u(\theta))\to 1/f_j$, and $\kappa_j(u(\theta))\ge 0.9/f_j$ for $\theta$ sufficiently close to $\pi/2$.

\paragraph{HL interpretation.}
This shows why P6/ledger invariants cannot replace HL-SPT-DELOC (and cannot enforce legality by themselves): if feasibility contains localized directions, energy conservation does not prevent legality loss.

\paragraph{Experiment pointer.}
NG-LEDGER is instantiated by an orthonormal two-channel feasible basis and a discrete rotation schedule; see \nolinkurl{paper/figures/energy_and_kappa_over_time.png}.

\subsection{NG-DIST: global distance-to-feasibility can be tiny while deep-band $\kappa_j$ is near maximal}

\begin{proposition}[NG-DIST: tiny global leakage can yield near-maximal deep-band concentration]\label{prop:ng-dist}
Fix a depth $j$.
For any target leakage ratio $r\in(0,1)$ and any target threshold fraction $\rho\in(0,1)$ (e.g.\ $\rho=0.9$), there exists a multiscale state $u\in H$ and a decomposition
\[
u = v + w,
\]
with the following properties:
\begin{enumerate}
\item (\emph{$v$ is SPT-legal}) for every depth $\ell$, $\kappa_\ell(v_\ell)\le 2$ (in particular $v$ is legal for $K_\kappa=2$).
\item (\emph{$w$ is globally tiny}) $\|w\|/\|v\| \le r$.
\item (\emph{$w$ is deep and localized}) $w_\ell=0$ for $\ell\neq j$, and $w_j$ is supported in a single cell of $\mathcal{C}_j$.
\item (\emph{Deep-band illegality}) the deep component satisfies
\[
\kappa_j(u_j)\ge \rho/f_j.
\]
\end{enumerate}
Thus, even when the global leakage ratio $\|w\|/\|v\|$ is arbitrarily small, the deep-band legality audit can be violated by a factor on the order of $1/f_j$.
\end{proposition}

\paragraph{Proof sketch (explicit scaling).}
Let $v=v^{\mathrm{coarse}}+v^{\mathrm{deep}}$ with $v^{\mathrm{coarse}}\in H_0$ and $v^{\mathrm{deep}}\in H_j$.
Choose $\|v^{\mathrm{coarse}}\|=1$ and let $v^{\mathrm{deep}}$ be a cell-uniform vector in $H_j$ of norm $\delta>0$, so $\kappa_j(v^{\mathrm{deep}})=1$ and hence $v$ is legal with threshold $2$ for $\delta$ sufficiently small.
Let $w$ be supported in a single cell at depth $j$ with $\|w\|=r$ and chosen orthogonal to $v^{\mathrm{deep}}$ (e.g.\ in a distinct internal coordinate of that cell).
Then $\|w\|/\|v\|\approx r$ because $\|v\|\approx 1$, but the deep band satisfies $u_j=v^{\mathrm{deep}}+w$ with $\|w\|\gg \|v^{\mathrm{deep}}\|$ if $\delta\ll r$.
In that case the maximizing cell contains almost all the deep-band energy, and one checks from Definitions~\ref{def:cell-conc}--\ref{def:kappa-def} that $\kappa_j(u_j)$ can be made $\ge \rho/f_j$ by choosing $\delta \le r/3$ (or smaller).
This is exactly the mechanism quantified by the cellwise defect control in HL-SPT-ROBUST$\infty$.

\paragraph{Experiment pointer.}
NG-DIST is instantiated by constructing a legal multiband baseline with a tiny deep-band component, then adding a cell-localized perturbation to that deep band; see \nolinkurl{paper/figures/leakage_vs_kappa.png}.

\subsection{FM-INV: leakage out of feasibility can cause instant illegality}

\begin{proposition}[FM-INV: without feasibility invariance, legality is not forward-closed]\label{prop:fm-inv}
Assume the artifact slots HL-SPT-DELOC, HL-SPT-FRAME, HL-SPT-MODE hold, so that $K^{\mathrm{pred}}_{\mathrm{glob}}$ in \eqref{eq:Kpred-global} is finite.
Suppose a dynamics admits a time $t_\star$ and a depth $j$ such that at time $t_\star$ the band component decomposes as
\[
u_j(t_\star)=v+w,
\qquad v\in\mathcal{F}_j,
\qquad w\perp \mathcal{F}_j,
\]
where $\|w\|/\|u(t_\star)\|\le 10^{-2}$ (small global leakage), but $w$ is supported in a single cell of $\mathcal{C}_j$ and dominates $v$ in that cell.
Then it is possible that
\[
\kappa_j(u_j(t_\star)) \;>\; K^{\mathrm{pred}}_{\mathrm{glob}},
\]
so SPT-legality with threshold $K^{\mathrm{pred}}_{\mathrm{glob}}$ fails at time $t_\star$ despite all artifact slots holding.
In particular, legality is not forward-closed without HL-SPT-INV or the robust substitute HL-SPT-ROBUST$\infty$.
\end{proposition}

\paragraph{Proof sketch.}
This is a dynamical packaging of NG-DIST (Prop.~\ref{prop:ng-dist}): choose $v$ feasible and (by T-SPT-00) legal within the predicted envelope, then inject a cell-localized off-feasible component $w$ that is globally tiny but deep-band dominant in one cell, making $\kappa_j(v+w)$ exceed the envelope.
Such an injection is exactly what HL-SPT-INV forbids, and what HL-SPT-ROBUST$\infty$ controls in a way that is aligned with the legality audit.

\paragraph{Experiment pointer.}
FM-INV is instantiated by injecting a localized off-feasible component into an otherwise invariant ladder evolution; see \nolinkurl{paper/figures/kappa_and_leakage_highlight.png}.
 \section{Toy models (non-PDE)}\label{sec:toy-models}

This section records simple toy dynamical systems---purely finite-dimensional---that satisfy (or violate) the HL slots.
They are designed to illustrate that the persistence theorems are about \emph{representation compatibility} rather than analytic smoothing.
The corresponding numerical implementations appear in the experiment suite (Section~12).

\subsection{Toy model A: feasible coefficient ladder (exact invariance)}

Fix depths $j=0,\dots,J_{\max}$ and channel maps $\Gamma_j:\mathbb{R}^{m_j}\to H_j$ with synthesis matrices $\Psi_j$ (columns are channels).
Consider the discrete-time coefficient ladder dynamics
\begin{equation}\label{eq:toy-ladder}
a_j(t+1)
\;=\;
A_j a_j(t)
\;+\;
B_{j\leftarrow j-1}a_{j-1}(t)
\;+\;
B_{j\leftarrow j+1}a_{j+1}(t)
\;+\;
\xi_j(t),
\end{equation}
where $A_j\in\mathbb{R}^{m_j\times m_j}$ is stable, $B_{j\leftarrow k}$ are coupling matrices of compatible sizes, and $\xi_j(t)\in\mathbb{R}^{m_j}$ is a coefficient-space noise/forcing term.
Define the band states by
\begin{equation}\label{eq:toy-ladder-state}
u_j(t) \;:=\; \Gamma_j a_j(t)\in H_j,
\qquad
u(t) \;:=\; \bigoplus_{j=0}^{J_{\max}} u_j(t)\in H.
\end{equation}

\begin{theorem}[Toy model A (invariant ladder) implies legality persistence]\label{thm:toy-ladder-invariant}
Assume the artifact slots HL-SPT-DELOC, HL-SPT-FRAME, and HL-SPT-MODE hold for the channel maps $\Gamma_j$ (Section~\ref{sec:hypotheses}), and assume the trajectory is initialized in the feasible sets: $u_j(0)\in\mathcal{F}_j$ for all $j$.
Then the ladder dynamics \eqref{eq:toy-ladder}--\eqref{eq:toy-ladder-state} satisfies exact feasibility invariance HL-SPT-INV:
\[
u_j(t)\in\mathcal{F}_j \quad \text{for all } t \text{ and all } j,
\]
and therefore SPT-legality persists for all time with threshold
\[
K_\kappa = K^{\mathrm{pred}}_{\mathrm{glob}}
\]
as in Theorem~\ref{thm:t-spt-01}.
\end{theorem}

\paragraph{Explanation.}
The invariance is immediate from the construction $u_j(t)=\Gamma_j a_j(t)$: all evolution occurs in coefficient space and is synthesized into $H_j$ by $\Gamma_j$.
The theorem then follows by applying T-SPT-01.

\paragraph{Experiment pointer.}
This model is instantiated in Exp-SPT-01-PERSIST; see \nolinkurl{paper/figures/kappa_heatmap.png} and \nolinkurl{paper/figures/kappa_max_over_time.png}.

\subsection{Toy model B: invariant ladder with localized off-feasible injection (break invariance)}

To isolate FM-INV, we perturb Toy model A by injecting an off-feasible component in a chosen depth $j_\star$ at some time $t_\star$.
Let $E_{j_\star}$ be the packaging projector onto $\mathcal{F}_{j_\star}$ and fix a vector $w_{\mathrm{inj}}\in H_{j_\star}$ with $(I-E_{j_\star})w_{\mathrm{inj}}=w_{\mathrm{inj}}$ that is supported in a single cell of $\mathcal{C}_{j_\star}$.
Define the injected dynamics at time $t_\star$ by
\begin{equation}\label{eq:toy-inj}
u_{j_\star}(t_\star^+) \;=\; u_{j_\star}(t_\star) + \eta\,w_{\mathrm{inj}},
\end{equation}
with injection strength $\eta>0$ and all other depths unchanged.

\begin{proposition}[Toy model B (leakage injection) can cause immediate illegality]\label{prop:toy-ladder-leakage}
Assume the artifact slots HL-SPT-DELOC, HL-SPT-FRAME, and HL-SPT-MODE hold so that $K^{\mathrm{pred}}_{\mathrm{glob}}$ is finite.
There exist choices of depth $j_\star$ (with small $f_{j_\star}$), injection time $t_\star$, and off-feasible cell-localized $w_{\mathrm{inj}}$ such that:
\begin{enumerate}
\item the global leakage ratio $\|\eta w_{\mathrm{inj}}\|/\|u(t_\star)\|$ can be made arbitrarily small, yet
\item the post-injection state violates legality:
\[
\kappa_{j_\star}(u_{j_\star}(t_\star^+)) \;>\; K^{\mathrm{pred}}_{\mathrm{glob}}.
\]
\end{enumerate}
Thus legality is not forward-closed without HL-SPT-INV (or without a robust substitute such as HL-SPT-ROBUST$\infty$), in agreement with Prop.~\ref{prop:fm-inv}.
\end{proposition}

\paragraph{Explanation.}
This is a dynamical realization of the NG-DIST mechanism: the injected component is globally tiny but concentrated in one deep cell, so it can dominate the deep-band max-cell fraction.
The injection violates HL-SPT-INV by construction.

\paragraph{Experiment pointer.}
This failure mode is instantiated in FM-INV; see \nolinkurl{paper/figures/kappa_and_leakage_highlight.png}.
 \section{Numerical experiments and assumption audits}\label{sec:experiments}

This section reports numerical experiments that harden the theorem overlays and, equally importantly, harden the \emph{sharpness boundary}.
All experiments are performed in the finite-dimensional Hilbert/band/cell model of Section~\ref{sec:definitions}.
They are not used to ``discover'' the theory; they are used to:
(i) verify that the happy-path hypotheses are sufficient in a mechanized setting, and
(ii) demonstrate that dropping any critical HL slot produces the intended failure mode.

\subsection{Protocol: standardized outputs, exports, and audits}

Each experiment run produces a standardized results bundle (metadata, diagnostics, summaries, timeseries, and figures) and an automated assumption audit report (\texttt{audit.json}).
For the present paper we freeze:
(i) a manifest-driven ``paper run set'' and
(ii) machine-readable exports collected across runs.

\paragraph{Frozen exports in the repository.}
The key quantitative numbers used for later tables are stored in:
\begin{itemize}
\item \texttt{paper/artifacts/key\_metrics.csv}: per-run scalar metrics (max observed $\kappa$, predicted envelopes, ratios, and highlight parameters).
\item \nolinkurl{paper/artifacts/audit_matrix.csv}: per-run HL slot flags (delocalization/\allowbreak frame/\allowbreak mode pass/fail and overall pass).
\item \texttt{paper/artifacts/figures\_index.csv}: inventory of run figures by run id.
\end{itemize}
These exports are produced automatically by the runner and collector scripts (see Appendix~A).

\paragraph{Replicates (not one-seed figures).}
To reduce the risk of ``one-seed miracles,'' we additionally ran a small replicate batch for the headline invariant persistence experiment (Exp-SPT-01-PERSIST) using five seeds (0--4) under the paper-scale configuration.
The resulting per-run summaries and aggregate statistics are frozen in \texttt{paper/artifacts/replicates/}, and Figure~\ref{fig:spt01-replicates-hist} shows the distribution of the maximum observed $\kappa$ across seeds.

\paragraph{How to read the plots.}
All plots are tied to explicit theorem constants.
In particular:
\begin{itemize}
\item The feasibility legality envelope $K^{\mathrm{feas}}_j$ and global envelope $K^{\mathrm{pred}}_{\mathrm{glob}}$ are defined in \eqref{eq:kfeas-def} and \eqref{eq:Kpred-global}.
\item The robust threshold is $K^{\mathrm{rob}}(\varepsilon_\infty)$ in \eqref{eq:t-spt-02-Krob}.
\item Failure mode plots are interpreted as intentional violations of specific HL slots (Section~\ref{sec:hypotheses}) and the corresponding no-go propositions in Section~\ref{sec:nogo}.
\end{itemize}

\begin{table}[t]
\centering
\small
\caption{Key numerical metrics for the frozen paper run set. Predicted bounds are taken from the theorem overlays (Definitions~\ref{def:kfeas} and \eqref{eq:Kpred-global}); pass/fail flags are from \texttt{paper/artifacts/audit\_matrix.csv}.}
\label{tab:key-metrics}
\setlength{\tabcolsep}{5pt}\begin{tabular}{@{}lrrrcccc@{}}
\toprule
Experiment & Max observed $\kappa$ & Predicted bound & Ratio & DELOC & FRAME & MODE & Overall \\
\midrule
\texttt{spt00\_bound} & 3.96 & 110 & 0.0359 & \(\checkmark\) & \(\checkmark\) & \(\checkmark\) & \(\checkmark\) \\
\texttt{spt01\_persist} & 3.22 & 10.2 & 0.316 & \(\checkmark\) & \(\checkmark\) & \(\checkmark\) & \(\checkmark\) \\
\texttt{spt02\_robust\_cell} & -- & -- & -- & \(\checkmark\) & \(\checkmark\) & \(\checkmark\) & \(\checkmark\) \\
\texttt{spt04\_fm\_inv} & 54.8 & 13 & 4.23 & \(\checkmark\) & \(\checkmark\) & \(\checkmark\) & \(\checkmark\) \\
\texttt{spt04\_fm\_deloc} & 256 & -- & -- & \(\times\) & \(\checkmark\) & \(\checkmark\) & \(\times\) \\
\texttt{spt04\_ng\_dist} & -- & -- & -- & \(\checkmark\) & \(\checkmark\) & \(\checkmark\) & \(\checkmark\) \\
\texttt{spt04\_ng\_ledger} & 256 & -- & -- & \(\times\) & \(\checkmark\) & \(\checkmark\) & \(\times\) \\
\bottomrule
\end{tabular}
\end{table}
 
\subsection{Happy path experiments: theorem regime hardening}

\subsubsection{Exp-SPT-00-BOUND: static feasibility $\Rightarrow$ legality}

Exp-SPT-00-BOUND samples coefficients $a\sim N(0,I)$, forms feasible vectors $v=\Gamma_j a$, and records $\kappa_j(v)$ across varying channel counts $m_j$.
The purpose is to validate the scaling predicted by T-SPT-00 (Theorem~\ref{thm:t-spt-00}): feasible concentration is controlled by a conservative envelope comparable to $K^{\mathrm{feas}}_j$.

\begin{figure}[t]
\centering
\includegraphics[width=0.92\linewidth]{figures/kappa\_vs\_m\_violin.png}
\caption{Exp-SPT-00-BOUND (static feasibility). Distributions of $\kappa_j(\Gamma_j a)$ across sampled coefficients and varying $m_j$, compared against the predicted feasibility envelope (T-SPT-00). Exact numeric maxima and ratios are recorded in \texttt{paper/artifacts/key\_metrics.csv}.}
\label{fig:exp-spt-00-bound}
\end{figure}

\subsubsection{Exp-SPT-01-PERSIST: exact invariance $\Rightarrow$ persistence}

Exp-SPT-01-PERSIST instantiates Toy model A (Section~\ref{sec:toy-models}): a coefficient-space ladder evolution $a_j(t)$ coupled across depths, with band states synthesized by $u_j(t)=\Gamma_j a_j(t)$.
By construction this satisfies HL-SPT-INV (exact feasibility invariance), so T-SPT-01 (Theorem~\ref{thm:t-spt-01}) predicts legality persistence for the global envelope $K^{\mathrm{pred}}_{\mathrm{glob}}$.

\begin{figure}[t]
\centering
\includegraphics[width=0.92\linewidth]{figures/kappa\_heatmap.png}
\caption{Exp-SPT-01-PERSIST (exact invariance). Heatmap of $\kappa_j(u_j(t))$ across time and depth for the invariant ladder system. In the happy-path regime, observed $\kappa$ remains bounded and consistent with the predicted global envelope; see \texttt{paper/artifacts/key\_metrics.csv}.}
\label{fig:exp-spt-01-persist}
\end{figure}

\begin{figure}[t]
\centering
\includegraphics[width=0.78\linewidth]{figures/spt01\_persist\_max\_kappa\_hist.png}
\caption{Replicate hardening for Exp-SPT-01-PERSIST (five seeds, paper-scale config). Histogram of the maximum observed $\kappa$ per run. Aggregate CSV summaries are frozen under \texttt{paper/artifacts/replicates/}.}
\label{fig:spt01-replicates-hist}
\end{figure}

\subsubsection{Exp-SPT-02-ROBUST-CELL: robust persistence under cellwise defect}

Exp-SPT-02-ROBUST-CELL adds a controlled off-feasible perturbation while enforcing a \emph{cellwise} defect constraint of the form
$\mathrm{def}^\infty_j(u_j(t))\le \varepsilon_\infty f_j$
(Definition~\ref{def:defect-infty}, HL-SPT-ROBUST$\infty$).
T-SPT-02 (Theorem~\ref{thm:t-spt-02}) predicts a degraded but explicit envelope $K^{\mathrm{rob}}(\varepsilon_\infty)$.

\begin{figure}[t]
\centering
\includegraphics[width=0.92\linewidth]{figures/max\_kappa\_vs\_epsilon.png}
\caption{Exp-SPT-02-ROBUST-CELL (robust invariance). Max observed $\kappa$ as a function of the imposed cellwise defect level $\varepsilon_\infty$, compared against the robust envelope $K^{\mathrm{rob}}(\varepsilon_\infty)$ from T-SPT-02.}
\label{fig:exp-spt-02-robust}
\end{figure}

\subsection{Sharpness experiments: break one HL slot at a time}

The remaining experiments are designed to be \emph{diagnostic counterexamples}.
Each one violates a specific HL slot and demonstrates the corresponding failure mode from Section~\ref{sec:nogo}.

\subsubsection{FM-INV: tiny leakage can cause a large $\kappa$ spike}

FM-INV injects a localized off-feasible component into an otherwise invariant ladder evolution (Toy model B, Prop.~\ref{prop:toy-ladder-leakage}).
This violates HL-SPT-INV and demonstrates Prop.~\ref{prop:fm-inv}: global leakage can be tiny while deep-band $\kappa$ immediately violates the predicted envelope.

\begin{figure}[t]
\centering
\includegraphics[width=0.92\linewidth]{figures/kappa\_and\_leakage\_highlight.png}
\caption{FM-INV (break invariance). A localized off-feasible injection produces an immediate deep-band $\kappa$ blow-up despite small global leakage, illustrating Prop.~\ref{prop:fm-inv} and motivating cellwise defect control (HL-SPT-ROBUST$\infty$).}
\label{fig:fm-inv}
\end{figure}

\subsubsection{FM-FRAME: delocalized columns without conditioning do not protect combinations}

FM-FRAME keeps per-column delocalization bounded while worsening the conditioning of $\Gamma_j$.
This isolates Prop.~\ref{prop:fm-frame}: delocalized columns do not prevent extreme concentration tails without HL-SPT-FRAME.

\begin{figure}[t]
\centering
\includegraphics[width=0.92\linewidth]{figures/kappa\_tail\_vs\_cond.png}
\caption{FM-FRAME (ill conditioning). Extreme concentration events grow rapidly as $\mathrm{cond}(\Gamma_j)$ worsens, even while per-column delocalization remains bounded. This illustrates Prop.~\ref{prop:fm-frame}.}
\label{fig:fm-frame}
\end{figure}

\subsubsection{FM-DELOC and FM-MODE: feasibility can contain or approximate wavepackets}

FM-DELOC includes an explicitly cell-supported channel.
This demonstrates Prop.~\ref{prop:fm-deloc}: feasibility can contain a wavepacket-like direction achieving $\kappa=1/f_j$.
FM-MODE, in contrast, illustrates Prop.~\ref{prop:fm-mode}: when $m_j$ grows too fast with depth, depth-uniform legality thresholds disappear even if channels are individually well-behaved.

\begin{figure}[p]
\centering
\begin{minipage}{0.49\linewidth}
\centering
\includegraphics[width=\linewidth]{figures/localized\_channel\_cell\_energy.png}
\end{minipage}
\hfill
\begin{minipage}{0.49\linewidth}
\centering
\includegraphics[width=\linewidth]{figures/max\_kappa\_vs\_depth\_by\_regime.png}
\end{minipage}
\caption{Left: FM-DELOC (localized feasible direction) achieves the ceiling $\kappa=1/f_j$, illustrating Prop.~\ref{prop:fm-deloc}. Right: FM-MODE (mode explosion) shows that a fixed legality threshold fails when interface size grows too fast with depth, illustrating Prop.~\ref{prop:fm-mode}.}
\label{fig:fm-deloc-mode}
\end{figure}

\subsection{No-go experiments: why global metrics and ledger invariants do not suffice}

\subsubsection{NG-DIST: global distance to feasibility is misleading}

NG-DIST demonstrates Prop.~\ref{prop:ng-dist}: global leakage ratios can be arbitrarily small while deep-band $\kappa$ becomes near-maximal, because deep bands may carry tiny global norm.
This is why robust persistence must control a cellwise defect aligned with the max-cell audit.

\subsubsection{NG-LEDGER: energy conservation does not prevent concentration blow-up}

NG-LEDGER demonstrates Prop.~\ref{prop:ng-ledger}: an energy-preserving evolution within a feasible subspace can rotate a state from delocalized to localized while keeping $\|u\|$ constant.

\begin{figure}[p]
\centering
\begin{minipage}{0.49\linewidth}
\centering
\includegraphics[width=\linewidth]{figures/leakage\_vs\_kappa.png}
\end{minipage}
\hfill
\begin{minipage}{0.49\linewidth}
\centering
\includegraphics[width=\linewidth]{figures/energy\_and\_kappa\_over\_time.png}
\end{minipage}
\caption{Left: NG-DIST: tiny global leakage can coincide with near-ceiling deep-band concentration, illustrating Prop.~\ref{prop:ng-dist}. Right: NG-LEDGER: energy remains flat while $\kappa$ rises to the ceiling due to feasible localized directions, illustrating Prop.~\ref{prop:ng-ledger}.}
\label{fig:ng-dist-ledger}
\end{figure}

\subsection{Audit interpretation}

Across all experiments, the assumption audit report (\texttt{audit.json}) plays the role of the P6 certificate: it checks whether the artifact slots (HL-SPT-DELOC/FRAME/MODE) pass numerical thresholds and whether invariance/defect evidence is consistent with the intended regime (happy path) or intended failure (FM/NG).
The matrix export \texttt{paper/artifacts/audit\_matrix.csv} is a convenient summary, and the scalar export \texttt{paper/artifacts/key\_metrics.csv} records the quantitative gaps between observed maxima and predicted envelopes.
 \section{Assumption audit}\label{sec:assumption-audit}

This section is a meta-level ``assumption audit'' in the sense of the Six-Birds primitives \cite{Tsiokos2026SixBirds}.
The goal is to clarify which hypotheses are purely structural (representation-level choices), which are the true hard slots needed for closure, and how the numerical audits relate to theorem conditions.

\subsection{Structural vs.\ hard slots}

We organize the HL slots from Section~\ref{sec:hypotheses} into two classes.

\paragraph{Structural representation choices (design slots).}
These are choices about how the multiscale state is represented and packaged.
They are not ``forced'' by accounting; they must be built into the representation.

\begin{itemize}
\item \textbf{Packaging/feasibility geometry (P5/P2):} choosing channel maps $\Gamma_j$ and feasible sets $\mathcal{F}_j=\Ran(\Gamma_j)$, together with packaging projectors $E_j$ (Definition~\ref{def:packaging}).
\item \textbf{Staging index (P4):} choosing a depth index $j$ and a cell partition $\mathcal{C}_j$ that defines what counts as localization.
\item \textbf{Mode-growth law (HL-SPT-MODE):} committing to an interface-size budget $m_j$ across depth.
\end{itemize}

These design slots are where the Six-Birds picture is most concrete: packaged ``objects'' are the fixed points of $E_j$, and gating is the restriction to $\mathcal{F}_j$.

\paragraph{Hard slots (closure-enabling constraints).}
These are the quantitatively checkable conditions without which legality is not forward-closed.

\begin{itemize}
\item \textbf{HL-SPT-DELOC:} structural anti-localization in the artifact itself. Without it, feasibility can contain wavepacket-like directions and legality fails even statically (FM-DELOC, Prop.~\ref{prop:fm-deloc}).
\item \textbf{HL-SPT-FRAME:} nondegenerate frame bounds. Without it, delocalized columns do not protect combinations and concentration tails can explode (FM-FRAME, Prop.~\ref{prop:fm-frame}).
\item \textbf{HL-SPT-INV / HL-SPT-ROBUST$\infty$:} compatibility between evolution and packaging. Without it, small leakage can cause immediate deep-band illegality (FM-INV, Prop.~\ref{prop:fm-inv}); and robust persistence requires defect control formulated at the same cell resolution as the legality audit (T-SPT-02 and NG-DIST).
\end{itemize}

The role of accounting/audits (P6) is to \emph{check} these hard slots and detect impending failure, not to replace them.

\subsection{How the automated audit maps to HL slots}

Each experiment run emits a machine-readable audit report (\texttt{audit.json}) and the paper freezes a cross-run audit summary:
\texttt{paper/artifacts/audit\_matrix.csv}.
The intent is to make the HL slots executable and regression-testable.

\paragraph{Artifact checks.}
The audit evaluates:
\begin{itemize}
\item \textbf{HL-SPT-DELOC} via the empirical constant $\kappa^{\mathrm{emp}}_0(j)$ (Definition~\ref{def:kappa0-emp}) and a threshold $K_{\mathrm{deloc}}$.
\item \textbf{HL-SPT-FRAME} via per-depth conditioning estimates (Definition~\ref{def:frame-bounds}) and a threshold $K_{\mathrm{cond}}$.
\item \textbf{HL-SPT-MODE} by recording $m_j$ and comparing to the configured mode budget.
\end{itemize}
These checks are ``artifact audits'' in the sense that they depend only on $\Gamma_j$ and the cell partition, not on the dynamics.

\paragraph{System compatibility checks.}
For experiments that claim invariance or robustness, the audit also records evidence aligned with the corresponding slot:
\begin{itemize}
\item \textbf{HL-SPT-INV} is evidenced in runs where the state is constructed as $u_j(t)=\Gamma_j a_j(t)$, making feasibility exact by construction.
\item \textbf{HL-SPT-ROBUST$\infty$} is evidenced in runs where the perturbation is engineered to satisfy the bound
$\mathrm{def}^\infty_j(u_j(t))\le \varepsilon_\infty f_j$
(Definition~\ref{def:defect-infty}).
\end{itemize}

\paragraph{Scalar exports for theorem overlays.}
The paper also freezes per-run scalar summaries in \texttt{paper/artifacts/key\_metrics.csv}, including observed maxima of $\kappa$, predicted envelopes (e.g.\ $K^{\mathrm{pred}}_{\mathrm{glob}}$), and max-to-predicted ratios.
These are used to tie plots and tables directly to theorem constants rather than to informal descriptions.

\subsection{Upgrading empirical diagnostics into theorem conditions}

A numerical audit is not a proof.
However, it can be made \emph{theorem-aligned} if the reported quantities correspond directly to explicit hypotheses and constants in the theorems.

\paragraph{What is already theorem-aligned.}
The key observation is that T-SPT-00--T-SPT-02 are stated entirely in terms of measurable constants:
$\kappa_0(j)$, $(\alpha_j,\beta_j)$ (or their condition proxy), $m_j$, and (for robustness) $\varepsilon_\infty$.
The experiment suite computes empirical versions of exactly these quantities and uses them to form predicted envelopes (Definition~\ref{def:kfeas} and \eqref{eq:Kpred-global}).

\paragraph{Where the gap remains.}
To turn diagnostics into proofs, one would need:
\begin{itemize}
\item \textbf{A certified artifact bound:} prove (or mechanize) that the intended construction of channels yields a uniform bound $\kappa^{\mathrm{emp}}_0(j)\le K_{\mathrm{deloc}}$ and a uniform frame lower bound $\alpha_j\ge \alpha_{\min}>0$ across depth.
\item \textbf{A system-level invariance proof:} for a given dynamics, prove HL-SPT-INV (exact invariance) or prove a bound of the form \eqref{eq:HL-ROBUST} (robust invariance) with an explicit $\varepsilon_\infty$.
\item \textbf{A link between route mismatch and defect:} when dynamics is only approximately compatible with packaging, one must bound the per-cell mismatch, not just a global norm mismatch, because legality is a max-cell audit.
\end{itemize}

\paragraph{Six-Birds interpretation.}
In the Six-Birds framing \cite{Tsiokos2026SixBirds}, the audit layer (P6) is functioning exactly as intended:
it provides a mechanized certificate that the HL slots are satisfied in the happy-path regime and that the correct slot is violated in each failure-mode regime.
When a slot fails, the correct response is not to reinterpret the audit as a substitute for structure, but to redesign packaging/gating (P5/P2) or rewrite the operator (P1) until compatibility is restored.

\subsection{Summary}

The persistence theorems show that legality is a forward-closed regime only under explicit structural hypotheses.
The numerical audits ensure these hypotheses are not treated implicitly: they are executable, testable, and regressible while the paper is being written.
 \section{Conclusion}\label{sec:conclusion}

We introduced a minimal, mechanizable setting in which a multiscale ``legality'' property can be defined and studied without PDE-specific regularization.
Legality is quantified by a max-cell concentration audit $\kappa_j$ (Definitions~\ref{def:cell-conc}--\ref{def:kappa-def}), and a state is SPT-legal if $\kappa_j$ is uniformly bounded across depths (Definition~\ref{def:spt-legal}).
The core question was: when is this legal regime \emph{forward-closed} under time evolution?

\subsection{What we proved}

Our main results are structural closure theorems.

\begin{itemize}
\item \textbf{Feasibility geometry controls legality (T-SPT-00).}
If each depth admits a channelized feasibility set $\mathcal{F}_j=\Ran(\Gamma_j)$ whose columns are cell-delocalized and whose frame bounds are nondegenerate, then feasible vectors obey an explicit legality envelope scaling like
$\kappa^{\mathrm{emp}}_0(j)\,\mathrm{cond}(\Gamma_j)\,m_j$
(Theorem~\ref{thm:t-spt-00}).

\item \textbf{Exact invariance yields persistence (T-SPT-01).}
If dynamics preserves feasibility bandwise (HL-SPT-INV), then legality propagates forward in time with threshold given by the predicted global envelope $K^{\mathrm{pred}}_{\mathrm{glob}}$
(Theorem~\ref{thm:t-spt-01}).

\item \textbf{Robust persistence requires cellwise defect control (T-SPT-02).}
If feasibility is only approximately preserved, legality can still be forward-closed provided the mismatch is controlled \emph{cellwise} in the same partition used to define legality (HL-SPT-ROBUST$\infty$).
This yields an explicit degraded envelope $K^{\mathrm{rob}}(\varepsilon_\infty)$
(Theorem~\ref{thm:t-spt-02}).
\end{itemize}

\subsection{What we ruled out}

We also proved sharp boundary statements showing why these hypotheses are not optional.

\begin{itemize}
\item If feasibility contains localized directions, legality can fail even statically (FM-DELOC, Prop.\ \ref{prop:fm-deloc}), and energy/\allowbreak ledger invariants do not prevent blow-up (NG-LEDGER, Prop.~\ref{prop:ng-ledger}).
\item Delocalized columns do not protect combinations without a frame lower bound (FM-FRAME, Prop.~\ref{prop:fm-frame}).
\item Global distance-to-feasibility is not a robustness notion for a max-cell audit (NG-DIST, Prop.~\ref{prop:ng-dist}).
\end{itemize}

These no-go results clarify the nature of ``legal regime closure'': anti-localization must be enforced structurally in the artifact, and compatibility between evolution and packaging must be stated at the resolution of the audit.

\subsection{Six-Birds interpretation}

The paper can be read as an instance of the Six-Birds primitives \cite{Tsiokos2026SixBirds}.
Packaging (P5) is realized as the bandwise projector $E_j$ onto feasible sets, gating (P2) is the invariance/defect compatibility that keeps evolution within (or near) those fixed points, staged interface growth (P4) is the mode-growth control on $m_j$, and audits (P6) are the explicit legality and artifact diagnostics emitted per run.
Route mismatch/holonomy (P3) appears not as a certificate but as a measurable defect: the off-feasible component $(I-E_j)u_j(t)$ and, in the robust theorem, its cellwise ratio.
Finally, when closure fails, the appropriate Six-Birds move is operator rewrite (P1): modify packaging/gating or the effective operator until invariance (or robust invariance) becomes provable and auditable.

\subsection{Future work: SPT on higher-order artifacts}

The present work treats $\Gamma_j$ (and hence $\mathcal{F}_j$ and $E_j$) as fixed artifacts.
A natural next step is to apply the SPT perspective to the artifact itself: treat the channel maps $\Gamma_j$ as emergent objects produced by repeated packaging and route-stabilization.

One toy direction is to formalize an ``artifact update'' map
\[
\Gamma_j \longmapsto \mathcal{U}_j(\Gamma_j; \text{packaged environment data})
\]
and to study conditions under which iterated packaging yields an attracting family of artifacts with:
(i) bounded per-column delocalization, (ii) nondegenerate frame lower bounds, and (iii) slow mode growth.
In Six-Birds terms, this would express how repeated P5 packaging and P3 route-stabilization can produce a fixed point in which anti-localization is \emph{emergent} rather than assumed.
Even in finite-dimensional toy models, identifying conditions under which such an artifact-level fixed point exists would connect legality persistence to the formation of the representation itself.

\subsection{Reproducibility}

All figures and tables in this paper are generated by a manifest-driven runner and collected into frozen exports under \texttt{paper/artifacts/} and \texttt{paper/figures/}.
Appendix~A lists the commands to reproduce the full paper run set and to validate the expected HL slot behavior.
 
\appendix
\appendix
\section{Reproducibility checklist and commands}\label{app:repro}

This appendix records the commands used to generate the numerical artifacts referenced in the paper and to validate the expected hypothesis-slot behavior.
The repository is organized so that the paper can be built without rerunning any experiments: canonical figures and exports are copied into \texttt{paper/figures/} and \texttt{paper/artifacts/} (see \texttt{paper/artifacts/ARTIFACTS.md} for the source-to-destination mapping).

\subsection{Paper build}

From the repository root, build the PDF using \texttt{latexmk}:
\begin{verbatim}
latexmk -pdf -interaction=nonstopmode -cd -output-directory=build paper/main.tex
\end{verbatim}

\subsection{Generate the full paper run set (manifest-driven)}

The ``paper run set'' is generated by running all YAML configs under \texttt{configs/paper/} and writing results into \texttt{results/\_paper/} with a manifest:

\begin{verbatim}
python -m spt.paper_runall --configs_dir configs/paper --out results/_paper
\end{verbatim}

The command fails fast if any experiment fails.
The repository also contains \texttt{configs/ci/} which are reduced-parameter analogues used for fast regression tests; the figures in this paper correspond to \texttt{configs/paper/}.

\subsection{Collect exports (tables and figure index)}

After generating the run set, collect machine-readable exports:

\begin{verbatim}
python -m spt.paper_collect --paper_dir results/_paper \
    --out results/_paper/_exports
\end{verbatim}

This produces, among others:
\texttt{key\_metrics.csv}, \texttt{audit\_matrix.csv}, and \texttt{figures\_index.csv}.
The paper freezes copies of these files under \texttt{paper/artifacts/} for self-contained builds.
Each run directory also stores provenance files including \texttt{resolved\_config.yaml} (the fully resolved configuration), \texttt{env.json} (environment snapshot), and \texttt{run\_metadata.json} (code version hash, timestamps, and seeds), enabling exact reruns from the recorded contract.

\subsection{Validate runs and enforce expected failure modes}

To validate the run set and check that intentional failure-mode experiments violate the correct slots:

\begin{verbatim}
python -m spt.validate_paper --paper_dir results/_paper
\end{verbatim}

For a single run directory:

\begin{verbatim}
python -m spt.validate_run results/_paper/<run_id>
\end{verbatim}

\subsection{Generate the LaTeX key-metrics table from frozen exports}

The small key-metrics table included in Section~\ref{sec:experiments} is generated from the frozen CSV exports:

\begin{verbatim}
python scripts/make_paper_tables.py
\end{verbatim}

This writes \nolinkurl{paper/tables/key_metrics.tex}, which is included by:
\begin{quote}
\verb|\begin{table}[t]
\centering
\small
\caption{Key numerical metrics for the frozen paper run set. Predicted bounds are taken from the theorem overlays (Definitions~\ref{def:kfeas} and \eqref{eq:Kpred-global}); pass/fail flags are from \texttt{paper/artifacts/audit\_matrix.csv}.}
\label{tab:key-metrics}
\setlength{\tabcolsep}{5pt}\begin{tabular}{@{}lrrrcccc@{}}
\toprule
Experiment & Max observed $\kappa$ & Predicted bound & Ratio & DELOC & FRAME & MODE & Overall \\
\midrule
\texttt{spt00\_bound} & 3.96 & 110 & 0.0359 & \(\checkmark\) & \(\checkmark\) & \(\checkmark\) & \(\checkmark\) \\
\texttt{spt01\_persist} & 3.22 & 10.2 & 0.316 & \(\checkmark\) & \(\checkmark\) & \(\checkmark\) & \(\checkmark\) \\
\texttt{spt02\_robust\_cell} & -- & -- & -- & \(\checkmark\) & \(\checkmark\) & \(\checkmark\) & \(\checkmark\) \\
\texttt{spt04\_fm\_inv} & 54.8 & 13 & 4.23 & \(\checkmark\) & \(\checkmark\) & \(\checkmark\) & \(\checkmark\) \\
\texttt{spt04\_fm\_deloc} & 256 & -- & -- & \(\times\) & \(\checkmark\) & \(\checkmark\) & \(\times\) \\
\texttt{spt04\_ng\_dist} & -- & -- & -- & \(\checkmark\) & \(\checkmark\) & \(\checkmark\) & \(\checkmark\) \\
\texttt{spt04\_ng\_ledger} & 256 & -- & -- & \(\times\) & \(\checkmark\) & \(\checkmark\) & \(\times\) \\
\bottomrule
\end{tabular}
\end{table}
 |
\end{quote}

\subsection{Where the frozen paper artifacts live}

For convenience and to support writing without rerunning experiments, the repository stores canonical paper artifacts in:

\begin{itemize}
\item \texttt{paper/figures/} (PNG figures referenced by the LaTeX source).
\item \texttt{paper/artifacts/} (CSV exports and \texttt{ARTIFACTS.md} mapping sources to destinations).
\item \texttt{paper/tables/} (generated \texttt{.tex} tables such as \texttt{key\_metrics.tex}).
\end{itemize}

\subsection{Optional: replicate batch for a headline experiment}

\begin{verbatim}
python -m spt.batch_run --config configs/paper/spt01_persist.yaml --seeds 0..4 \
  --aggregate_dir results/_aggregate/paper_spt01_persist_seeds0_4
\end{verbatim}

\subsection{Lean4 mechanization scaffolding (optional)}

A Lean4 + mathlib project is initialized in the repository for mechanization scaffolding.
The files under \texttt{Lean/SPT/} provide definition stubs (finite cell index sets, concentration, $\kappa$, legality, channel maps, feasible sets) and theorem statement stubs corresponding to the main theorems.
The Lean statements are currently placeholders (\texttt{sorry}) and are intended as a starting point for future mechanization rather than as part of the present proof.

To build the Lean project (from repo root):

\begin{verbatim}
lake build
\end{verbatim}
 
\section*{Acknowledgments}
Placeholder acknowledgments.

\bibliographystyle{plain}
\bibliography{refs}

\end{document}
